% ====================================================================
%  徐森林、薛春华《数学分析》（第一册）
%  第 5 章  不定积分
%  源：原书印刷页 282-327（PDF 物理页 = 印刷页 + 6）
% ====================================================================

\chapter{不定积分}\label{ch:5}

我们已经学会了求导运算, 并且大家都能熟练地、正确地计算初等函数的导数. 同时还学会了分段函数的求导, 特别在分界点处应按导数的定义来求得该点处的导数. 本章将介绍求导的逆运算, 也就是求 $f(x)$ 的原函数 $F(x)$, 这里 $F'(x)=f(x)$. 然后, 在区间上求得所有的原函数, 即不定积分 $\int f(x)\mathrm{d}x=F(x)+C$ ($C$ 为任意常数).

\section{原函数、不定积分}\label{sec:5.1}

\begin{definition}\label{def:5.1.1}
设 $f$ 为区间 $I$ 上的可导函数, 则 $\mathrm{d}f(x)=f'(x)\mathrm{d}x$ 称为 $f(x)$ 的微分.
\end{definition}

显然, 有下面定理.

\begin{theorem}\label{thm:5.1.1}
(1) 设 $f$, $g$ 在区间 $I$ 上可导, 则
\begin{gather*}
\mathrm{d}(f(x)\pm g(x))=\mathrm{d}f(x)\pm\mathrm{d}g(x);\\
\mathrm{d}(f(x)g(x))=f(x)\mathrm{d}g(x)+g(x)\mathrm{d}f(x),\quad
\mathrm{d}(cf(x))=c\,\mathrm{d}f(x)\quad (c\in\mathbb{R});\\
\mathrm{d}\left(\frac{f(x)}{g(x)}\right)=\frac{g(x)\mathrm{d}f(x)-f(x)\mathrm{d}g(x)}{g^{2}(x)}
\quad (g(x)\neq 0,\ \forall x\in I).
\end{gather*}
(2) 设 $f$ 在区间 $I$ 上可导, $g$ 在区间 $J$ 上可导, 且 $I\supseteq g(J)$. $y=f(u)$, $u=g(x)$, $y=(f\circ g)(x)$. 则
\[
\mathrm{d}\left[(f\circ g)(x)\right]=(f\circ g)'(x)\mathrm{d}x=f'(u)\mathrm{d}u
\]
或
\[
\mathrm{d}y=y_{x}'\mathrm{d}x=y_{u}'\mathrm{d}u \quad(\text{称它为 1 阶微分的不变性}),
\]
其中 $x$ 为自变量, $u$ 为中间变量. 此式表明, 无论 $u$ 为自变量还是中间变量, 其 1 阶微分形式是不变的.
\end{theorem}

\begin{proof}
(1) 以 $\frac{f(x)}{g(x)}$ 的微分为例.
\begin{align*}
\mathrm{d}\left(\frac{f(x)}{g(x)}\right)
&=\left(\frac{f(x)}{g(x)}\right)'\mathrm{d}x
=\frac{g(x)f'(x)-f(x)g'(x)}{g^{2}(x)}\mathrm{d}x\\
&=\frac{g(x)\mathrm{d}f(x)-f(x)\mathrm{d}g(x)}{g^{2}(x)}.
\end{align*}
其他各式可类似证明.

(2)
\[
\mathrm{d}\left[(f\circ g)(x)\right]=(f\circ g)'(x)\mathrm{d}x=f'\left(g(x)\right)g'(x)\mathrm{d}x=f'(u)\mathrm{d}u.
\]
\end{proof}

\begin{example}\label{ex:5.1.1}
设 $y=\mathrm{e}^{x^{2}+x}$, 求 $\mathrm{d}y$.
\end{example}

\begin{solution}[解法 1]
$\mathrm{d}y=(\mathrm{e}^{x^{2}+x})'\mathrm{d}x=\mathrm{e}^{x^{2}+x}(2x+1)\mathrm{d}x$.
\end{solution}

\begin{solution}[解法 2]
令 $u=x^{2}+x$, 则 $y=\mathrm{e}^{u}$,
\[
\mathrm{d}y=(\mathrm{e}^{u})'\mathrm{d}u=\mathrm{e}^{u}\mathrm{d}u
=\mathrm{e}^{x^{2}+x}\mathrm{d}(x^{2}+x)=\mathrm{e}^{x^{2}+x}(2x+1)\mathrm{d}x.
\]
\end{solution}

\begin{example}\label{ex:5.1.2}
设 $x^{3}+y^{3}=4$, 求 $\mathrm{d}y$.
\end{example}

\begin{solution}
对方程 $x^{3}+y^{3}=4$ 两边求微分 ($x$ 与 $y$ 地位平等) 得
\begin{gather*}
3x^{2}\mathrm{d}x+3y^{2}\mathrm{d}y=0,\\
y^{2}\mathrm{d}y=-x^{2}\mathrm{d}x,\\
y_{x}'=\frac{\mathrm{d}y}{\mathrm{d}x}=-\frac{x^{2}}{y^{2}},\quad
x_{y}'=\frac{\mathrm{d}x}{\mathrm{d}y}=-\frac{y^{2}}{x^{2}}.
\end{gather*}
或者解出 $y=\sqrt[3]{4-x^{3}}$, 求微分得
\begin{align*}
\mathrm{d}y&=\frac{1}{3}(4-x^{2})^{-\frac{2}{3}}(0-3x^{2})\mathrm{d}x=-\frac{x^{2}}{(4-x^{2})^{\frac{2}{3}}}\mathrm{d}x\\
&=-\frac{x^{2}}{y^{2}}\mathrm{d}x.
\end{align*}
\end{solution}

\begin{example}\label{ex:5.1.3}
求 $\arctan\frac{y}{x}=\ln\sqrt{x^{2}+y^{2}}$ 关于 $\mathrm{d}x$ 与 $\mathrm{d}y$ 的微分方程, 并求 $y'$.
\end{example}

\begin{solution}
对原式两边微分得
\begin{gather*}
\frac{\mathrm{d}\left(\frac{y}{x}\right)}{1+\left(\frac{y}{x}\right)^{2}}
=\frac{1}{2}\frac{\mathrm{d}(x^{2}+y^{2})}{x^{2}+y^{2}},\\
\frac{x^{2}\frac{x\,\mathrm{d}y-y\,\mathrm{d}x}{x^{2}}}{x^{2}+y^{2}}
=\frac{1}{2}\cdot\frac{2(x\,\mathrm{d}x+y\,\mathrm{d}y)}{x^{2}+y^{2}},\\
x\,\mathrm{d}y-y\,\mathrm{d}x=x\,\mathrm{d}x+y\,\mathrm{d}y,\\
(x-y)\mathrm{d}y=(x+y)\mathrm{d}x.
\end{gather*}
于是
\[
y'=\frac{\mathrm{d}y}{\mathrm{d}x}=\frac{x+y}{x-y}.
\]
\end{solution}

\begin{definition}\label{def:5.1.2}
在区间 $I$ 上, 如果 $f$ 为 $F$ 的导函数, 即 $F'(x)=f(x)$ 或 $\mathrm{d}F=F'(x)\mathrm{d}x=f(x)\mathrm{d}x$, 则称 $F$ 为 $f$ 在区间 $I$ 上的一个原函数.
\end{definition}

\begin{theorem}\label{thm:5.1.2}
设 $f(x)$ 在区间 $I$ 上有原函数 $F(x)$, 即在 $I$ 上 $F'(x)=f(x)$, 则函数族 $\{F(x)+C\mid C\in\mathbb{R}\}$ 为 $f(x)$ 的全部原函数.
\end{theorem}

\begin{proof}
在区间 $I$ 上, 由推论~\ref{cor:3.3.2} 知,
\begin{align*}
G(x)\ \text{为}\ f(x)\ \text{的原函数}
&\Longleftrightarrow G'(x)=f(x)\\
&\Longleftrightarrow G'(x)=F'(x)\\
&\Longleftrightarrow G(x)=F(x)+C,\ \text{其中}\ C\ \text{为常数}.
\end{align*}
\end{proof}

\begin{definition}\label{def:5.1.3}
设函数 $f(x)$ 在区间上有一个原函数 $F(x)$, 则函数族 $\{F(x)+C\}$ 就称为 $f(x)$ 在 $I$ 上的不定积分, 记作
\[
\int f(x)\mathrm{d}x=F(x)+C,
\]
其中 $C$ 称为不定常数, $\int$ 称为积分号, $f(x)$ 称为被积函数, $x$ 称为积分变量, $f(x)\mathrm{d}x$ 称为被积表达式或被积形式.
\end{definition}

由不定积分的定义立即有下面结论:

(1)
\[
\left(\int f(x)\mathrm{d}x\right)'=\left(F(x)+C\right)'=f(x),
\]
或
\[
\mathrm{d}\left(\int f(x)\mathrm{d}x\right)=\mathrm{d}(F(x)+C)=F'(x)\mathrm{d}x=f(x)\mathrm{d}x;
\]

(2)
\[
\int F'(x)\mathrm{d}x=F(x)+C,
\]
或
\[
\int\mathrm{d}F(x)=F(x)+C.
\]

\begin{theorem}\label{thm:5.1.3}
设函数 $f(x)$ 与 $g(x)$ 在区间 $I$ 上都存在原函数, 对 $\forall k_{1},k_{2}\in\mathbb{R}$, 且 $k_{1}^{2}+k_{2}^{2}\neq 0$, 则 $k_{1}f(x)+k_{2}g(x)$ 在 $I$ 上也存在原函数, 且
\[
\int\left[k_{1}f(x)+k_{2}g(x)\right]\mathrm{d}x
=k_{1}\int f(x)\mathrm{d}x+k_{2}\int g(x)\mathrm{d}x.
\]
\end{theorem}

\begin{proof}
设 $F(x)$ 与 $G(x)$ 分别为 $f(x)$ 与 $g(x)$ 的一个原函数, 即 $F'(x)=f(x)$, $G'(x)=g(x)$. 于是
\[
\left[k_{1}F(x)+k_{2}G(x)\right]'=k_{1}F'(x)+k_{2}G'(x)=k_{1}f(x)+k_{2}g(x),
\]
即 $k_{1}F(x)+k_{2}G(x)$ 为 $k_{1}f(x)+k_{2}g(x)$ 的一个原函数. 由此得到
\begin{align*}
\int\left[k_{1}f(x)+k_{2}g(x)\right]\mathrm{d}x
&=k_{1}F(x)+k_{2}G(x)+C\\
&=k_{1}\left[F(x)+C_{1}\right]+k_{2}\left[G(x)+C_{2}\right]\\
&=k_{1}\int f(x)\mathrm{d}x+k_{2}\int g(x)\mathrm{d}x,
\end{align*}
其中 $C=k_{1}C_{1}+k_{2}C_{2}$.
\end{proof}

现在我们来考虑不定积分的几何意义. 如果 $F(x)$ 为 $f(x)$ 在区间 $I$ 上的一个原函数, 那么称 $y=F(x)$ 的图形为 $f(x)$ 的一条积分曲线. 于是, $f(x)$ 的不定积分 $\int f(x)\mathrm{d}x=F(x)+C$ 在几何上表示由 $f(x)$ 的某条积分曲线 $y=F(x)$ 沿纵轴方向任意平移所得的切积分曲线组成的曲线族. 显然, 在每一条积分曲线上横坐标相同的点处做切线, 则这些切线互相平行, 其斜率均为 $[F(x)+C]'=F'(x)=f(x)$ (图~\ref{fig:5.1.1}).

\begin{figure}[htbp]
\centering
\begin{tikzpicture}[x=1.5cm,y=1.05cm,>=Stealth]
  % 坐标轴
  \draw[->] (-0.3,0) -- (6.2,0) node[below] {$x$};
  \draw[->] (0,-1.9) -- (0,2.1) node[left] {$y$};
  \node[below left] at (0,0) {$O$};

  % 三条平移曲线（同一形状，沿 y 轴平移）
  \draw[thick,domain=0.2:5.6,samples=100,smooth]
    plot (\x, {1.05 + 0.42*sin(deg(\x*0.95)) - 0.16*cos(deg(\x*0.7))});
  \draw[thick,domain=0.2:5.6,samples=100,smooth]
    plot (\x, {0.32 + 0.42*sin(deg(\x*0.95)) - 0.16*cos(deg(\x*0.7))});
  \draw[thick,domain=0.2:5.6,samples=100,smooth]
    plot (\x, {-0.75 + 0.42*sin(deg(\x*0.95)) - 0.16*cos(deg(\x*0.7))});

  % x0 竖线
  \draw[dashed] (3.0,0) -- (3.0,-1.15);
  \node[below] at (3.0,0) {$x_{0}$};

  % 曲线标注
  \node[right] at (5.6, 1.62) {$y=F(x)+C$};
  \node[right] at (5.6, 0.86) {$y=F(x)$};
  \node[right] at (4.7, -0.22) {$y=[F(x)-F(x_{0})]+y_{0}$};
\end{tikzpicture}
\caption{}\label{fig:5.1.1}
\end{figure}

在求原函数的具体问题中, 往往先求出全体原函数 $F(x)+C$, 然后从中确定一个满足 $F(x_{0})+C=y_{0}$ (称为初始条件, 由具体问题所确定) 的原函数, 它对应着积分曲线族中通过点 $(x_{0},y_{0})=(x_{0},F(x_{0})+C)$ 的那一条积分曲线:
\begin{gather*}
F(x_{0})+C=y_{0},\quad C=y_{0}-F(x_{0}),\\
F(x)+C=F(x)+y_{0}-F(x_{0})=\left[F(x)-F(x_{0})\right]+y_{0}.
\end{gather*}

下一章将给出的微积分基本定理~\ref{thm:6.3.2} 指出连续函数 $f(x)$ 必有一个原函数 $\int_{x_{0}}^{x}f(t)\mathrm{d}t$, 因此
\[
\int_{x_{0}}^{x}f(t)\mathrm{d}t+y_{0}
\]
为满足初始条件的原函数, 即它所代表的曲线经过点 $(x_{0},y_{0})$.

再考虑不定积分的物理意义. 设质点作直线运动, 离原点的位移为 $x=x(t)$, 瞬时速度为 $v(t)=x'(t)$, 加速度为 $a(t)=v'(t)=x''(t)$. 反过来, 已知加速度如何求速度, 再求位移, 这就要用到求原函数与不定积分. 具体表达为
\[
v(t)=\int a(t)\mathrm{d}t,\qquad
x(t)=\int v(t)\mathrm{d}t=\int\left[\int a(t)\mathrm{d}t\right]\mathrm{d}t.
\]
例如: 质点作匀加速直线运动, 它的加速度为 $a(t)=v'(t)=x''(t)=a$ (常数), 初始位移为 $x_{0}=x(t_{0})$, 初始速度 $v_{0}=v(t_{0})=x'(t_{0})$. 于是,
\begin{align*}
v(t)&=\int a\mathrm{d}t=at+C_{1}=at+(v_{0}-at_{0})=a(t-t_{0})+v_{0},\\
x(t)&=\int v(t)\mathrm{d}t=\int\left[a(t-t_{0})+v_{0}\right]\mathrm{d}t\\
&=\frac{a}{2}(t-t_{0})^{2}+v_{0}t+C_{2}\\
&=\frac{a}{2}(t-t_{0})^{2}+v_{0}t+(x_{0}-v_{0}t_{0}).
\end{align*}

此外, 我们在下章注~\ref{rem:6.3.1} 中还可应用定积分的方法来表达上述 $v(t)$ 与 $x(t)$. 由求导公式, 立即得到如下基本积分公式表 (其中 $C,C_{1},C_{2}$ 为常数):

(1) $\int 0\,\mathrm{d}x=C$.

(2)
\[
\int x^{\alpha}\mathrm{d}x=\frac{x^{\alpha+1}}{\alpha+1}+C\quad (\alpha\neq -1,\ x>0),
\]
特别地
\[
\int x^{n}\mathrm{d}x=\frac{x^{n+1}}{n+1}+C,\quad x\in(-\infty,+\infty).
\]

(3) $\int\frac{\mathrm{d}x}{x}=\ln|x|+C$ ($x>0$ 或 $x<0$),
\[
\int\frac{\mathrm{d}x}{x}=
\begin{cases}
\ln x+C_{1}, & x>0,\\
\ln(-x)+C_{2}, & x<0.
\end{cases}
\]

(4)
\[
\int a^{x}\mathrm{d}x=\frac{a^{x}}{\ln a}+C\quad (a>0,\ a\neq 1),\qquad
\int\mathrm{e}^{x}\mathrm{d}x=\mathrm{e}^{x}+C.
\]

(5)
\[
\int\sin x\,\mathrm{d}x=-\cos x+C,\qquad
\int\cos x\,\mathrm{d}x=\sin x+C.
\]

(6)
\begin{align*}
\int\sec^{2}x\,\mathrm{d}x&=\int\frac{\mathrm{d}x}{\cos^{2}x}=\tan x+C,\\
\int\csc^{2}x\,\mathrm{d}x&=\int\frac{\mathrm{d}x}{\sin^{2}x}=-\cot x+C,\\
\int\sec x\tan x\,\mathrm{d}x&=\sec x+C,\qquad
\int\csc x\cot x\,\mathrm{d}x=-\csc x+C.
\end{align*}

(7)
\[
\int\frac{\mathrm{d}x}{1+x^{2}}=\arctan x+C_{1}=-\arctan x+C_{2}.
\]

(8)
\[
\int\frac{\mathrm{d}x}{\sqrt{1-x^{2}}}=\arcsin x+C_{1}=-\arccos x+C_{2}.
\]

(9)
\[
\int\sinh x\,\mathrm{d}x=\cosh x+C,\qquad
\int\cosh x\,\mathrm{d}x=\sinh x+C.
\]
其中 $\sinh x=\frac{\mathrm{e}^{x}-\mathrm{e}^{-x}}{2}$, $\cosh x=\frac{\mathrm{e}^{x}+\mathrm{e}^{-x}}{2}$.

另外由
\[
\left(\ln(x+\sqrt{a^{2}+x^{2}})\right)'=\frac{1+\frac{2x}{2\sqrt{a^{2}+x^{2}}}}{x+\sqrt{a^{2}+x^{2}}}=\frac{1}{\sqrt{a^{2}+x^{2}}}
\]
可得

(10) $\int\frac{\mathrm{d}x}{\sqrt{a^{2}+x^{2}}}=\ln(x+\sqrt{a^{2}+x^{2}})+C$.

从定理~\ref{thm:5.1.3} 及上面不定积分的一些简单性质可得到许多简单函数的原函数.

\begin{example}\label{ex:5.1.4}
(1)
\[
\int(a_{0}x^{n}+a_{1}x^{n-1}+\cdots+a_{n-1}x+a_{n})\mathrm{d}x
=\frac{a_{0}x^{n+1}}{n+1}+\frac{a_{1}x^{n}}{n}+\cdots+\frac{a_{n-1}x^{2}}{2}+a_{n}x+C.
\]
(2)
\[
\int\frac{x^{2}}{1+x^{2}}\mathrm{d}x=\int\left(1-\frac{1}{1+x^{2}}\right)\mathrm{d}x=x-\arctan x+C.
\]
(3)
\[
\int\left(3x^{2}+\frac{4}{x}\right)\mathrm{d}x=x^{3}+4\ln|x|+C\quad (x>0\ \text{或}\ x<0).
\]
(4)
\begin{align*}
\int\frac{x^{2}+1}{\sqrt{x}}\mathrm{d}x&=\int\left(x^{\frac{3}{2}}+x^{-\frac{1}{2}}\right)\mathrm{d}x
=\frac{x^{\frac{5}{2}}}{\frac{5}{2}}+\frac{x^{\frac{1}{2}}}{\frac{1}{2}}+C\\
&=\frac{2x^{\frac{5}{2}}}{5}+2x^{\frac{1}{2}}+C.
\end{align*}
(5)
\begin{align*}
\int\frac{x^{4}+1}{x^{2}+1}\mathrm{d}x&=\int\left(x^{2}-1+\frac{2}{x^{2}+1}\right)\mathrm{d}x\\
&=\frac{x^{3}}{3}-x+2\arctan x+C.
\end{align*}
(6)
\[
\int\frac{3^{x+3}}{2^{x}}\mathrm{d}x=27\int\left(\frac{3}{2}\right)^{x}\mathrm{d}x
=27\frac{\left(\frac{3}{2}\right)^{x}}{\ln\frac{3}{2}}+C.
\]
(7)
\begin{align*}
\int(10^{x}-10^{-x})^{2}\mathrm{d}x&=\int(10^{2x}+10^{-2x}-2)\mathrm{d}x\\
&=\frac{(10^{2})^{x}}{\ln 10^{2}}+\frac{(10^{-2})^{x}}{\ln 10^{-2}}-2x+C\\
&=\frac{1}{2\ln 10}(10^{2x}-10^{-2x})-2x+C.
\end{align*}
(8)
\[
\int\tan^{2}x\,\mathrm{d}x=\int(\sec^{2}x-1)\mathrm{d}x=\tan x-x+C.
\]
(9)
\begin{align*}
\int\frac{\cos 2x}{\cos x+\sin x}\mathrm{d}x&=\int\frac{\cos^{2}x-\sin^{2}x}{\cos x+\sin x}\mathrm{d}x\\
&=\int(\cos x-\sin x)\mathrm{d}x=\sin x+\cos x+C.
\end{align*}
(10)
\begin{align*}
\int\frac{\mathrm{d}x}{\cos^{2}x\sin^{2}x}&=\int\frac{\cos^{2}x+\sin^{2}x}{\cos^{2}x\sin^{2}x}\mathrm{d}x\\
&=\int(\csc^{2}x+\sec^{2}x)\mathrm{d}x=-\cot x+\tan x+C.
\end{align*}
(11)
\begin{align*}
\int\cos 3x\sin x\,\mathrm{d}x&=\frac{1}{2}\int(\sin 4x-\sin 2x)\mathrm{d}x\\
&=\frac{1}{2}\left(-\frac{\cos 4x}{4}+\frac{1}{2}\cos 2x\right)+C\\
&=-\frac{1}{8}\cos 4x+\frac{1}{4}\cos 2x+C.
\end{align*}
\end{example}

\begin{example}\label{ex:5.1.5}
已知曲线 $y=F(x)$ 的切线斜率为 $2x$, 求此曲线族. 如果曲线又经过点 $(0,1)$, 此曲线如何?
\end{example}

\begin{solution}
因为 $F'(x)=2x$, 所以
\[
F(x)=\int 2x\,\mathrm{d}x=x^{2}+C.
\]
又因为曲线经过点 $(0,1)$, 故
\[
1=F(0)=0^{2}+C=C,
\]
从而
\[
F(x)=x^{2}+1.
\]
见图~\ref{fig:5.1.2}. 或者应用下章定积分得到
\[
F(x)=\int_{0}^{x}2x\,\mathrm{d}x+F(0)=x^{2}+1.
\]
\end{solution}

\begin{figure}[htbp]
\centering
\begin{tikzpicture}[x=1.5cm,y=0.55cm,>=Stealth]
  \draw[->] (-2.3,0) -- (2.3,0) node[below] {$x$};
  \draw[->] (0,-0.8) -- (0,5.2) node[left] {$y$};
  \node[below left] at (0,0) {$O$};

  \draw[thick,domain=-2.15:2.15,samples=80,smooth] plot (\x, {(\x)*(\x)+0.6});
  \draw[thick,domain=-2.0:2.0,samples=80,smooth] plot (\x, {(\x)*(\x)+1.8});
  \draw[thick] (0,0.6) -- (0,1.8);
  \fill (0,1.8) circle (1.6pt);

  \node[right] at (1.55, 4.1) {$y=F(x)=x^{2}+1$};
  \node[right] at (1.55, 2.35) {$y=x^{2}+C$};
\end{tikzpicture}
\caption{}\label{fig:5.1.2}
\end{figure}

\begin{example}\label{ex:5.1.6}
求不定积分 $\int|x-1|\,\mathrm{d}x$.
\end{example}

\begin{solution}
设
\[
f(x)=|x-1|=
\begin{cases}
x-1, & x\geqslant 1,\\
1-x, & x<1.
\end{cases}
\]
显然它为连续函数. 根据第 6 章微积分基本定理~\ref{thm:6.3.2} 可知, $f(x)=|x-1|$ 有原函数 $F(x)$, 即 $F'(x)=f(x)$, 由此可看出 $F(x)$ 连续. 因为
\[
\int|x-1|\,\mathrm{d}x=
\begin{cases}
\displaystyle\int(x-1)\mathrm{d}x=\frac{x^{2}}{2}-x+C_{1}, & x>0,\\[2mm]
\displaystyle\int(1-x)\mathrm{d}x=x-\frac{x^{2}}{2}+C_{2}, & x<0.
\end{cases}
\]
令 $C_{1}=0$, 则得 $f(x)=|x-1|$ 的一个原函数为
\[
F(x)=
\begin{cases}
\frac{x^{2}}{2}-x, & x>1,\\[2mm]
x-\frac{x^{2}}{2}+C_{2}, & x<0.
\end{cases}
\]
由于 $F(x)$ 连续, 当然在 $x=1$ 处连续, 即
\[
-\frac{1}{2}=\lim_{x\to1^{+}}\left(\frac{x^{2}}{2}-x\right)
=\lim_{x\to1^{-}}\left(x-\frac{x^{2}}{2}+C_{2}\right)=\frac{1}{2}+C_{2},
\]
解得 $C_{2}=-1$, 代入上式, 有
\[
F(x)=
\begin{cases}
\frac{x^{2}}{2}-x, & x\geqslant 1,\\[2mm]
x-\frac{x^{2}}{2}-1, & x<1.
\end{cases}
\]
见图~\ref{fig:5.1.3}. 从而
\[
\int|x-1|\,\mathrm{d}x=F(x)+C.
\]
\end{solution}

\begin{figure}[htbp]
\centering
\begin{tikzpicture}[x=1.6cm,y=1.5cm,>=Stealth]
  \draw[->] (-1.6,0) -- (3.2,0) node[below] {$x$};
  \draw[->] (0,-1.6) -- (0,2.1) node[left] {$y$};
  \node[below left] at (0,0) {$O$};

  \draw[dashed] (1,0) -- (1,-0.5);
  \node at (1,0.12) {\scriptsize $1$};
  \node[below] at (2,0) {\scriptsize $2$};
  \draw[dashed] (0,-0.5) -- (1.15,-0.5);
  \node[left] at (0,-0.5) {\scriptsize $-\frac{1}{2}$};

  % y = x^2/2 - x (x>=1)
  \draw[very thick,domain=1:2.6,samples=50,smooth] plot (\x, {(\x)*(\x)/2-(\x)});
  % y = x - x^2/2 - 1 (x<1)
  \draw[very thick,domain=-1.4:1,samples=50,smooth] plot (\x, {(\x)-(\x)*(\x)/2-1});

  \node[right] at (2.2,1.35) {$y=\frac{x^{2}}{2}-x$};
  \node[below] at (2.2,0.62) {$y=F(x)$};
  \node[left] at (-1.0,-1.1) {$y=x-\frac{x^{2}}{2}-1$};
\end{tikzpicture}
\caption{}\label{fig:5.1.3}
\end{figure}

\begin{example}\label{ex:5.1.7}
求不定积分
\[
\int|\sin x|\,\mathrm{d}x.
\]
\end{example}

\begin{solution}
因为 $|\sin x|$ 在 $\mathbb{R}$ 上连续, 所以根据第 6 章微积分基本定理~\ref{thm:6.3.2} 知, $|\sin x|$ 有原函数 $F(x)$, 即 $F'(x)=|\sin x|$. 显然 $F'(x)\geqslant 0$ 且
\[
\{x\mid F'(x)=|\sin x|=0\}=\{k\pi\mid k\in\mathbb{Z}\}
\]
不构成区间, 故 $F(x)$ 为严格增的连续函数. 令
\[
F(x)=
\begin{cases}
-\cos x+c_{k}, & 2k\pi\leqslant x<(2k+1)\pi,\\
\cos x+d_{k}, & (2k+1)\pi\leqslant x<(2k+2)\pi.
\end{cases}
\]
取 $c_{0}=0$, 由 $F(x)$ 连续推得
\[
1+c_{k}=-\cos(2k+1)\pi+c_{k}=\cos(2k+1)\pi+d_{k}=-1+d_{k},
\]
所以
\[
d_{k}=c_{k}+2.
\]
再由
\[
1+d_{k}=\cos(2k+2)\pi+d_{k}=-\cos(2k+2)\pi+c_{k+1}=-1+c_{k+1},
\]
得到
\[
c_{k+1}=d_{k}+2.
\]
综合上述有 $c_{k}=4k$, $d_{k}=4k+2$. 从而
\[
F(x)=
\begin{cases}
-\cos x+4k, & 2k\pi\leqslant x<(2k+1)\pi,\\
\cos x+4k+2, & (2k+1)\pi\leqslant x<(2k+2)\pi,\ k\in\mathbb{Z}.
\end{cases}
\]
如图~\ref{fig:5.1.4} 所示. 而所要求的不定积分为
\[
\int|\sin x|\,\mathrm{d}x=F(x)+C.
\]
\end{solution}

\begin{figure}[htbp]
\centering
\begin{tikzpicture}[x=1.2cm,y=0.72cm,>=Stealth]
  \draw[->] (-4.1,0) -- (4.1,0) node[below] {$x$};
  \draw[->] (0,-3.3) -- (0,1.9) node[left] {$y$};
  \node[below left] at (0,0) {$O$};

  % y = sin x（原书用细虚线）
  \draw[dashed,domain=-3.5:3.5,samples=140,smooth] plot (\x, {sin(deg(\x))});
  \node[right] at (3.55,-1.45) {$y=\sin x$};

  % F(x)：严格增的分段函数，在 kπ 处相切连接
  \draw[very thick,domain=0:3.14159,samples=60,smooth] plot (\x, {-cos(deg(\x))});
  \draw[very thick,domain=-3.14159:0,samples=60,smooth] plot (\x, {cos(deg(\x))-2});
  \node[above] at (1.9,1.15) {$y=F(x)$};

  \draw (-1.5708,0.07) -- (-1.5708,-0.07);
  \node[below] at (-1.5708,-0.05) {\scriptsize $-\frac{\pi}{2}$};
  \draw (1.5708,0.07) -- (1.5708,-0.07);
  \node[below] at (1.5708,-0.05) {\scriptsize $\frac{\pi}{2}$};
\end{tikzpicture}
\caption{}\label{fig:5.1.4}
\end{figure}

\begin{example}\label{ex:5.1.8}
设
\[
F(x)=
\begin{cases}
x^{2}\sin\frac{1}{x}, & x\neq 0,\\
0, & x=0,
\end{cases}
\]
求 $F'(x)$.
\end{example}

\begin{solution}
\[
f(x)=F'(x)=
\begin{cases}
2x\sin\frac{1}{x}-\cos\frac{1}{x}, & x\neq 0,\\
0, & x=0.
\end{cases}
\]
显然, 它只有一个不连续点 $0$, 但 $f(x)$ 有原函数 $F(x)$.
\end{solution}

\begin{example}\label{ex:5.1.9}
设
\[
F(x)=
\begin{cases}
x^{\frac{3}{2}}\sin\frac{1}{x}, & x\neq 0,\\
0, & x=0,
\end{cases}
\]
求 $F'(x)$.
\end{example}

\begin{solution}
\[
f(x)=F'(x)=
\begin{cases}
\frac{3}{2}x^{\frac{1}{2}}\sin\frac{1}{x}-\frac{1}{x^{\frac{1}{2}}}\cos\frac{1}{x}, & x\neq 0,\\
0, & x=0.
\end{cases}
\]
\end{solution}

\begin{remark}\label{rmk:5.1.1}
显然, $f(x)$ 在 $x=0$ 附近无界, 根据定理~\ref{thm:6.1.2}, 它在含 $0$ 的任何闭区间上都非 Riemann 可积. 因此, 不能应用推论~\ref{cor:6.3.1} 由变上限积分得到 $f(x)$ 的原函数. 但 $f(x)$ 确有原函数 $F(x)$.
\end{remark}

\begin{example}\label{ex:5.1.10}
任何函数 $f$, 是否一定有原函数?
\end{example}

\begin{solution}
不一定. 反例 1: Dirichlet 函数
\[
D(x)=
\begin{cases}
1, & x\in\mathbb{Q},\\
0, & x\in\mathbb{R}\setminus\mathbb{Q}
\end{cases}
\]
无原函数.

证明参看例~\ref{ex:3.3.11}.

反例 2: 含第 1 类间断点的函数一定无原函数. 如
\[
\operatorname{sgn}x=
\begin{cases}
1, & x>0,\\
0, & x=0,\\
-1, & x<0
\end{cases}
\]
以 $0$ 为第 1 类间断点.

可用反证法证明. 假设 $f(x)$ 有原函数 $F(x)$, 即 $F'(x)=f(x)$. 由例~\ref{ex:3.3.12} 知, $f(x)=F'(x)$ 无第 1 类间断点, 这与 $f(x)$ 含第 1 类间断点相矛盾.
\end{solution}

如果一个关于 $x$ 与 $y$ 的微分等式通过分离变量后可化为
\[
f(y)\mathrm{d}y=g(x)\mathrm{d}x,
\]
则两边积分得到
\[
\int f(y)\mathrm{d}y=\int g(x)\mathrm{d}x.
\]
称上述方法为分离变量法.

\begin{example}\label{ex:5.1.11}
从物理实验知, 放射性元素的衰变速度与当时的质量成正比, 求衰变规律及半生期 $T$ (质量衰变到一半时的时间称为该放射性元素的半生期).
\end{example}

\begin{solution}
设放射性元素在时刻 $t$ 的质量为 $m=m(t)$, 则
\[
\frac{\mathrm{d}m}{\mathrm{d}t}=m'(t)=-km,
\]
其中负号表示质量减少. 经分离变量后得到
\begin{gather*}
\frac{\mathrm{d}m}{m}=-k\,\mathrm{d}t,\\
\int\frac{\mathrm{d}m}{m}=-k\int\mathrm{d}t,\\
\ln|m|=-kt+C_{1},\\
|m|=\mathrm{e}^{-kt+C_{1}}=\mathrm{e}^{C_{1}}\mathrm{e}^{-kt}=C_{2}\mathrm{e}^{-kt},\\
m=m(t)=C\mathrm{e}^{-kt}.
\end{gather*}
当 $t=0$ 时, 质量为 $m(0)$, 则 $m(0)=C\mathrm{e}^{-k\cdot 0}=C$,
\[
m=m(t)=m(0)\mathrm{e}^{-kt}.
\]
考虑该放射性元素的半生期 $T$, 它满足
\begin{gather*}
m(0)\mathrm{e}^{-kT}=\frac{1}{2}m(0),\\
\mathrm{e}^{-kT}=\frac{1}{2},\\
-kT=\ln\frac{1}{2}=-\ln 2,
\end{gather*}
\[
T=\frac{\ln 2}{k}.
\]
例如, 镭的半生期 $T=1600$ 年, 就是说 $1\,\mathrm{g}$ 镭经过 $1600$ 年后剩下 $0.5\,\mathrm{g}$.

此外, 我们在第 6 章注~\ref{rem:6.3.1} 中还可应用定积分的方法来表达上述 $m(t)$.
\end{solution}

\begin{example}\label{ex:5.1.12}
设 $f'(\arctan x)=x^{2}$, 求 $f(x)$.
\end{example}

\begin{solution}
设 $\arctan x=t$, 则 $x=\tan t$, $f'(t)=\tan^{2}t$,
\begin{align*}
f(t)&=\int\tan^{2}t\,\mathrm{d}t=\int(\sec^{2}t-1)\mathrm{d}t=\tan t-t+C,\\
f(x)&=\tan x-x+C.
\end{align*}
\end{solution}

\begin{exblock}{练习题 5.1}
  \item 求不定积分:
  \begin{enumerate}[label=(\arabic*)]
    \item $\int(2+x^{5})^{2}\mathrm{d}x$;
    \item $\int\left(\frac{1-x}{x}\right)^{2}\mathrm{d}x$;
    \item $\int\left(1-\frac{1}{x^{2}}\right)\sqrt{x}\,\mathrm{d}x$;
    \item $\int\frac{x^{5}}{1+x}\mathrm{d}x$;
    \item $\int\frac{x^{4}}{1+x^{2}}\mathrm{d}x$;
    \item $\int\left(1-x+x^{3}-\frac{1}{\sqrt[3]{x^{2}}}\right)\mathrm{d}x$;
    \item $\int\left(x-\frac{1}{\sqrt{x}}\right)^{2}\mathrm{d}x$;
    \item $\int\frac{x^{2}}{3(1+x^{2})}\mathrm{d}x$;
    \item $\int\sqrt{x\sqrt{x\sqrt{x}}}\,\mathrm{d}x$;
    \item $\int\left(\sqrt{\frac{1+x}{1-x}}+\sqrt{\frac{1-x}{1+x}}\right)\mathrm{d}x$;
    \item $\int(2^{x}+3^{x})^{2}\mathrm{d}x$;
    \item $\int 10^{t}\cdot 3^{2t}\,\mathrm{d}t$;
    \item $\int(\mathrm{e}^{x}-\mathrm{e}^{-x})^{3}\mathrm{d}x$;
    \item $\int\frac{\mathrm{e}^{3x}+1}{\mathrm{e}^{x}+1}\mathrm{d}x$;
    \item $\int\frac{\mathrm{d}x}{(x+a)(x+b)}$;
    \item $\int\frac{3}{\sqrt{4-4x^{2}}}\mathrm{d}x$;
    \item $\int\cosh x\,\mathrm{d}x$;
    \item $\int\sinh x\,\mathrm{d}x$;
    \item $\int\cos^{2}x\,\mathrm{d}x$;
    \item $\int\sin^{2}x\,\mathrm{d}x$;
    \item $\int\sqrt{1-\sin 2x}\,\mathrm{d}x$;
    \item $\int\frac{\cos 2x}{\cos x-\sin x}\mathrm{d}x$;
    \item $\int\frac{\cos 2x}{\cos^{2}x\sin^{2}x}\mathrm{d}x$;
    \item $\int(\cos x+\sin x)^{2}\mathrm{d}x$;
    \item $\int\tan^{2}x\,\mathrm{d}x$;
    \item $\int\cos x\cos 2x\,\mathrm{d}x$;
    \item $\int\frac{2^{x+1}-5^{x-1}}{10^{x}}\mathrm{d}x$;
    \item $\int|x|\,\mathrm{d}x$;
    \item $\int\mathrm{e}^{-|x|}\,\mathrm{d}x$;
    \item $\int\frac{x^{2}}{(1-x)^{100}}\mathrm{d}x$;
    \item $\int\frac{\sqrt{x}-2\sqrt[3]{x}-1}{\sqrt[4]{x}}\mathrm{d}x$.
  \end{enumerate}

  \item 求一曲线 $y=f(x)$, 使得在曲线上每一点 $(x,y)$ 处的切线斜率为 $3x^{2}$, 且通过点 $(2,9)$.

  \item 验证: $y=\frac{x^{2}}{2}\operatorname{sgn}x$ 为 $|x|$ 在 $(-\infty,+\infty)$ 上的一个原函数.

  \item 周期函数的原函数是否还是周期函数? 考察 $f(x)=|\sin x|$. 设 $f(x)$ 是一个周期为 $T$ 的函数, $F(x)$ 为它的一个原函数. 证明: $F(x)$ 有周期 $T\Longleftrightarrow F(T)=F(0)$. (提示: 考虑函数 $F(x+T)-F(x)$.)

  \item 设 $f(x)$ 为奇函数, $F(x)$ 为它的一个原函数. 证明: $F(x)$ 为偶函数. (提示: 考虑函数 $F(-x)-F(x)$.)

  \item 设 $f(x)$ 为偶函数, $F(x)$ 为它的一个原函数. 证明: $F(x)$ 为奇函数 $\Longleftrightarrow F(0)=0$. (提示: 考虑函数 $F(-x)+F(x)$.)
\end{exblock}

\section{换元积分法、分部积分法}\label{sec:5.2}

利用积分公式与积分的简单性质, 虽然能求出许多简单函数的原函数, 但是, 仅这一点点知识是很不够的, 还必须学会更多的计算原函数的方法和技巧, 换元 (变量代换) 积分法与分部积分法是两种重要的方法.

\begin{theorem}[换元 (变量代换) 积分法]\label{thm:5.2.1}
设 $g(u)$ 在 $[\alpha,\beta]$ 上有定义, $u=\varphi(x)$ 在 $[a,b]$ 上可导, 且 $\alpha\leqslant\varphi(x)\leqslant\beta$, $x\in[a,b]$. 并记
\[
f(x)=g(\varphi(x))\varphi'(x),\quad x\in[a,b].
\]
(1) 如果 $g(u)$ 在 $[\alpha,\beta]$ 上存在原函数 $G(u)$, 则 $f(x)$ 在 $[a,b]$ 上也存在原函数 $F(x)=G(\varphi(x))$. 于是
\begin{align*}
\int f(x)\mathrm{d}x&\xlongequal{\text{凑}}\int g(\varphi(x))\varphi'(x)\mathrm{d}x
\xlongequal{u=\varphi(x)}\int g(u)\mathrm{d}u\\
&=G(u)+C=G(\varphi(x))+C.
\end{align*}
(2) 如果 $\varphi'(x)\neq 0$, $x\in[a,b]$, 则上述可逆, 即当 $f(x)$ 在 $[a,b]$ 上存在原函数 $F(x)$ 时, $g(u)$ 在 $[\alpha,\beta]$ 上也存在原函数 $G(u)$, 且
\[
G(u)=F(\varphi^{-1}(u))+C,
\]
即
\begin{align*}
\int g(u)\mathrm{d}u&\xlongequal{u=\varphi(x)}\int g(\varphi(x))\varphi'(x)\mathrm{d}x\\
&=\int f(x)\mathrm{d}x=F(x)+C\\
&\xlongequal{x=\varphi^{-1}(u)}F(\varphi^{-1}(u))+C.
\end{align*}
\end{theorem}

\begin{proof}
(1) 由复合函数求导法知,
\begin{align*}
\left[G(\varphi(x))\right]'&=G'(\varphi(x))\varphi'(x)\\
&=g(\varphi(x))\varphi'(x)=f(x),
\end{align*}
所以, $G(\varphi(x))$ 为 $f(x)$ 的一个原函数.

(2) 在 $\varphi'(x)\neq 0$ 的条件下, $u=\varphi(x)$ 有反函数 $x=\varphi^{-1}(u)$, 且
\[
x'(u)=\frac{1}{\varphi'(x)}.
\]
于是, 有
\begin{align*}
\left[F(\varphi^{-1}(u))\right]'
&\xlongequal{x=\varphi^{-1}(u)}F'(x)x'(u)=f(x)\frac{1}{\varphi'(x)}\\
&=g(\varphi(x))\varphi'(x)\frac{1}{\varphi'(x)}\\
&=g(\varphi(x))=g(u),
\end{align*}
即 $F(\varphi^{-1}(u))$ 为 $g(u)$ 的一个原函数.
\end{proof}

换元 (变量代换) 积分法对应着复合函数的导数公式.

定理~\ref{thm:5.2.1}(1) 阐述了通过凑的方法可将一个不易积的不定积分 $\int f(x)\mathrm{d}x$ 化为一个较易积的不定积分 $\int g(\varphi(x))\varphi'(x)\mathrm{d}x=\int g(u)\mathrm{d}u$.

定理~\ref{thm:5.2.1}(2) 则通过主动做一个变量代换 $u=\varphi(x)$, 将不易积的不定积分 $\int g(u)\mathrm{d}u$ 化为一个较易积的不定积分 $\int g(\varphi(x))\varphi'(x)\mathrm{d}x=\int f(x)\mathrm{d}x$. 换元 (变量代换) 关键在于变量代换 $u=\varphi(x)$.

\begin{theorem}[分部积分法]\label{thm:5.2.2}
设 $u(x)$ 与 $v(x)$ 可导,\linebreak[4]不定积分 $\int u'(x)v(x)\mathrm{d}x$ 存在, 则 $\int u(x)v'(x)\mathrm{d}x$ 也存在, 且有
\[
\int u(x)v'(x)\mathrm{d}x=u(x)v(x)-\int u'(x)v(x)\mathrm{d}x,
\]
也可写作
\[
\int u(x)\mathrm{d}v(x)=u(x)v(x)-\int v(x)\mathrm{d}u(x).
\]
\end{theorem}

\begin{proof}
由乘积的求导公式得到
\begin{gather*}
[u(x)v(x)]'=u'(x)v(x)+u(x)v'(x),\\
u(x)v'(x)=[u(x)v(x)]'-u'(x)v(x).
\end{gather*}
由此立知, 不定积分 $\int u'(x)v(x)\mathrm{d}x$ 存在, 则 $\int u(x)v'(x)\mathrm{d}x$ 也存在, 且有
\begin{align*}
\int u(x)v'(x)\mathrm{d}x&=\int\left\{[u(x)v(x)]'-u'(x)v(x)\right\}\mathrm{d}x\\
&=u(x)v(x)-\int u'(x)v(x)\mathrm{d}x.
\end{align*}
\end{proof}

分部积分法对应着 $u$ 与 $v$ 之积 $uv$ 的导数公式 $(uv)'=u'v+uv'$ (导性).

从定理~\ref{thm:5.2.2} 可以看到, 一个不易积的不定积分 $\int uv'(x)\mathrm{d}x$ 通过分部积分法, 分出一部分 $uv$, 化为一个较易积的不定积分 $\int u'(x)v(x)\mathrm{d}x$.

\begin{example}\label{ex:5.2.1}
应用换元 (变量代换) 积分法求下列不定积分:
\begin{enumerate}[label=(\arabic*)]
  \item $\int\tan x\,\mathrm{d}x$;
  \item $\int x\mathrm{e}^{x^{2}}\mathrm{d}x$;
  \item $\int\frac{\mathrm{d}x}{ax+b}$ ($a\neq 0$);
  \item $\int\frac{\mathrm{d}x}{a^{2}+x^{2}}$ ($a>0$);
  \item $\int\frac{\mathrm{d}x}{x^{2}-a^{2}}$ ($a>0$);
  \item $\int\frac{x}{x^{2}-a^{2}}\mathrm{d}x$;
  \item $\int\frac{x}{x^{2}+a^{2}}\mathrm{d}x$;
  \item $\int\frac{\mathrm{d}x}{\sqrt{a^{2}-x^{2}}}$ ($a>0$);
  \item $\int\frac{\ln x}{x}\mathrm{d}x$;
  \item $\int\frac{\arcsin x}{\sqrt{1-x^{2}}}\mathrm{d}x$;
  \item $\int\cos^{2}x\,\mathrm{d}x$.
\end{enumerate}
\end{example}

\begin{solution}
(1)
\begin{align*}
\int\tan x\,\mathrm{d}x&=\int\frac{\sin x}{\cos x}\mathrm{d}x=-\int\frac{\mathrm{d}\cos x}{\cos x}
\xlongequal{u=\cos x}-\int\frac{\mathrm{d}u}{u}\\
&=-\ln|u|+C\xlongequal{u=\cos x}-\ln|\cos x|+C.
\end{align*}
(2)
\begin{align*}
\int x\mathrm{e}^{x^{2}}\mathrm{d}x&=\frac{1}{2}\int\mathrm{e}^{x^{2}}\mathrm{d}x^{2}
\xlongequal{t=x^{2}}\frac{1}{2}\int\mathrm{e}^{t}\mathrm{d}t\\
&=\frac{1}{2}\mathrm{e}^{t}+C=\frac{1}{2}\mathrm{e}^{x^{2}}+C.
\end{align*}
(3)
\begin{align*}
\int\frac{\mathrm{d}x}{ax+b}&=\frac{1}{a}\int\frac{\mathrm{d}(ax+b)}{ax+b}
\xlongequal{t=ax+b}\frac{1}{a}\int\frac{\mathrm{d}t}{t}\\
&=\frac{1}{a}\ln|t|+C=\frac{1}{a}\ln|ax+b|+C.
\end{align*}
(4)
\[
\int\frac{\mathrm{d}x}{a^{2}+x^{2}}=\frac{1}{a}\int\frac{\mathrm{d}\left(\frac{x}{a}\right)}{1+\left(\frac{x}{a}\right)^{2}}
=\frac{1}{a}\arctan\frac{x}{a}+C.
\]
(5)
\begin{align*}
\int\frac{\mathrm{d}x}{x^{2}-a^{2}}&=\frac{1}{2a}\int\left(\frac{1}{x-a}-\frac{1}{x+a}\right)\mathrm{d}x\\
&=\frac{1}{2a}\left[\int\frac{\mathrm{d}(x-a)}{x-a}-\int\frac{\mathrm{d}(x+a)}{x+a}\right]\\
&=\frac{1}{2a}\left[\ln|x-a|-\ln|x+a|\right]+C\\
&=\frac{1}{2a}\ln\left|\frac{x-a}{x+a}\right|+C.
\end{align*}
(6)
\[
\int\frac{x}{x^{2}-a^{2}}\mathrm{d}x=\frac{1}{2}\int\frac{\mathrm{d}(x^{2}-a^{2})}{x^{2}-a^{2}}
=\frac{1}{2}\ln|x^{2}-a^{2}|+C.
\]
(7)
\[
\int\frac{x}{x^{2}+a^{2}}\mathrm{d}x=\frac{1}{2}\int\frac{\mathrm{d}(x^{2}+a^{2})}{x^{2}+a^{2}}
=\frac{1}{2}\ln(x^{2}+a^{2})+C.
\]
(8)
\[
\int\frac{\mathrm{d}x}{\sqrt{a^{2}-x^{2}}}=\int\frac{\mathrm{d}\left(\frac{x}{a}\right)}{\sqrt{1-\left(\frac{x}{a}\right)^{2}}}
=\arcsin\frac{x}{a}+C\quad (a>0).
\]
(9)
\[
\int\frac{\ln x}{x}\mathrm{d}x=\int\ln x\,\mathrm{d}\ln x=\frac{\ln^{2}x}{2}+C.
\]
(10)
\[
\int\frac{\arcsin x}{\sqrt{1-x^{2}}}\mathrm{d}x=\int\arcsin x\,\mathrm{d}\arcsin x=\frac{(\arcsin x)^{2}}{2}+C.
\]
(11)
\begin{align*}
\int\cos^{2}x\,\mathrm{d}x&=\frac{1}{2}\int(\cos 2x+1)\mathrm{d}x
=\frac{1}{4}\int\cos 2x\,\mathrm{d}(2x)+\frac{1}{2}\int\mathrm{d}x\\
&=\frac{1}{4}\sin 2x+\frac{x}{2}+C.
\end{align*}
\end{solution}

\begin{example}\label{ex:5.2.2}
求 $\int\sec x\,\mathrm{d}x$.
\end{example}

\begin{solution}[解法 1]
\begin{align*}
\int\sec x\,\mathrm{d}x&=\int\frac{\mathrm{d}x}{\cos x}=\int\frac{\cos x}{\cos^{2}x}\mathrm{d}x
=\int\frac{\mathrm{d}\sin x}{1-\sin^{2}x}\\
&\xlongequal{\text{例}~\ref{ex:5.2.1}(5)}\frac{1}{2}\ln\frac{1+\sin x}{1-\sin x}+C.
\end{align*}
或进一步地, 有
\begin{align*}
\frac{1}{2}\ln\frac{1+\sin x}{1-\sin x}+C&=\frac{1}{2}\ln\frac{(1+\sin x)^{2}}{1-\sin^{2}x}+C\\
&=\frac{1}{2}\ln\left(\frac{1+\sin x}{\cos x}\right)^{2}+C=\ln\left|\frac{1+\sin x}{\cos x}\right|+C\\
&=\ln|\sec x+\tan x|+C.
\end{align*}
\end{solution}

\begin{solution}[解法 2]
\begin{align*}
\int\sec x\,\mathrm{d}x&=\int\frac{\sec x(\sec x+\tan x)}{\sec x+\tan x}\mathrm{d}x
=\int\frac{\mathrm{d}(\sec x+\tan x)}{\sec x+\tan x}\\
&=\ln|\sec x+\tan x|+C.
\end{align*}
\end{solution}

\begin{solution}[解法 3]
\begin{align*}
\int\sec x\,\mathrm{d}x&=\int\frac{\mathrm{d}x}{\cos x}=\int\frac{\mathrm{d}x}{\cos^{2}\frac{x}{2}-\sin^{2}\frac{x}{2}}\\
&=\int\frac{\sec^{2}\frac{x}{2}}{1-\tan^{2}\frac{x}{2}}\mathrm{d}x
=2\int\frac{\mathrm{d}\tan\frac{x}{2}}{1-\tan^{2}\frac{x}{2}}\\
&\xlongequal[\text{例}~\ref{ex:5.2.1}(5)]{t=\tan\frac{x}{2}}\ln\left|\frac{1+\tan\frac{x}{2}}{1-\tan\frac{x}{2}}\right|+C.
\end{align*}
或进一步地, 有
\begin{align*}
\ln\left|\frac{1+\tan\frac{x}{2}}{1-\tan\frac{x}{2}}\right|+C
&=\ln\left|\frac{\cos\frac{x}{2}+\sin\frac{x}{2}}{\cos\frac{x}{2}-\sin\frac{x}{2}}\right|+C
=\ln\left|\frac{\left(\cos\frac{x}{2}+\sin\frac{x}{2}\right)^{2}}{\cos^{2}\frac{x}{2}-\sin^{2}\frac{x}{2}}\right|+C\\
&=\ln\left|\frac{1+\sin x}{\cos x}\right|+C=\ln|\sec x+\tan x|+C.
\end{align*}
\end{solution}

同样方法可求
\begin{align*}
\int\csc x\,\mathrm{d}x&=-\ln\left|\frac{1+\cos x}{\sin x}\right|+C\\
&=-\ln|\csc x+\cot x|+C=-\ln\left|\cot\frac{x}{2}\right|+C.\quad (\text{留作习题.})
\end{align*}

\begin{example}\label{ex:5.2.3}
求 $\int\frac{\mathrm{d}x}{\sqrt{x^{2}-a^{2}}}$.
\end{example}

\begin{solution}[解法 1]
如图~\ref{fig:5.2.1} 所示.
\begin{align*}
\int\frac{\mathrm{d}x}{\sqrt{x^{2}-a^{2}}}
&\xlongequal{x=a\sec t}\int\frac{a\sec t\tan t}{a\tan t}\mathrm{d}t
=\int\sec t\,\mathrm{d}t\\
&\xlongequal{\text{例}~\ref{ex:5.2.2}}\ln|\sec t+\tan t|+C_{1}\\
&=\ln\left|\frac{x}{a}+\frac{\sqrt{x^{2}-a^{2}}}{a}\right|+C_{1}
=\ln\left|x+\sqrt{x^{2}-a^{2}}\right|+C.
\end{align*}
\end{solution}

\begin{solution}[解法 2]
当 $x>a$ 时, 令 $x=a\cosh t$, $t>0$, 则
\begin{align*}
\int\frac{\mathrm{d}x}{\sqrt{x^{2}-a^{2}}}
&\xlongequal{x=a\cosh t}\int\frac{a\sinh t}{a\sinh t}\mathrm{d}t\\
&=\int\mathrm{d}t=t+C_{1}\\
&=\ln\left|\frac{x}{a}+\sqrt{\left(\frac{x}{a}\right)^{2}-1}\right|+C_{1}\\
&=\ln\left|x+\sqrt{x^{2}-a^{2}}\right|+C,
\end{align*}
其中
\begin{gather*}
x=a\cosh t=a\frac{\mathrm{e}^{t}+\mathrm{e}^{-t}}{2},\\
(\mathrm{e}^{t})^{2}-2\frac{x}{a}\mathrm{e}^{t}+1=0,\\
\mathrm{e}^{t}=\frac{2\frac{x}{a}+2\sqrt{\left(\frac{x}{a}\right)^{2}-1}}{2}=\frac{x}{a}+\sqrt{\left(\frac{x}{a}\right)^{2}-1},\\
t=\ln\left|\frac{x}{a}+\sqrt{\left(\frac{x}{a}\right)^{2}-1}\right|.
\end{gather*}
显然, 对 $x<-a$, 上述结果也是成立的.
\end{solution}

\begin{figure}[htbp]
\centering
\begin{tikzpicture}[x=1.1cm,y=1.1cm,>=Stealth]
  \draw[thick] (0,0) -- (2.4,0) node[midway,below] {$\sqrt{x^{2}-a^{2}}$};
  \draw[thick] (0,0) -- (2.4,2.0);
  \draw[thick] (2.4,0) -- (2.4,2.0) node[midway,right] {$x$};
  \draw (2.15,0) -- (2.15,0.25) -- (2.4,0.25);
  \draw (0.5,0) arc (0:39.8:0.5);
  \node at (0.62,0.18) {$t$};
  \node[above] at (1.2,1.05) {$a$};
  \node[right] at (3.0,1.7) {$\sec t=\frac{x}{a}$};
  \node[right] at (3.0,1.1) {$\tan t=\frac{\sqrt{x^{2}-a^{2}}}{a}$};
\end{tikzpicture}
\caption{}\label{fig:5.2.1}
\end{figure}

做变量代换 $x=a\sec t$ 与 $x=a\cosh t$ 都是为了将 $\sqrt{x^{2}-a^{2}}$ 中的根号去掉, 然后化为一个易积的不定积分或已有结果的不定积分.

\begin{example}\label{ex:5.2.4}
求 $\int\frac{\mathrm{d}x}{\sqrt{x^{2}+a^{2}}}$ ($a>0$).
\end{example}

\begin{solution}[解法 1]
读者一定还记得例~\ref{ex:3.1.6}(2) 中, 有
\[
\left(\ln(x+\sqrt{x^{2}+a^{2}})\right)'=\frac{1}{\sqrt{x^{2}+a^{2}}},
\]
故
\[
\int\frac{\mathrm{d}x}{\sqrt{x^{2}+a^{2}}}=\ln(x+\sqrt{x^{2}+a^{2}})+C.
\]
\end{solution}

\begin{solution}[解法 2]
如图~\ref{fig:5.2.2} 所示. 如果事先不知道解法 1 中的导数公式, 可以进行变量代换 $x=a\tan t$ 得到
\begin{align*}
\int\frac{\mathrm{d}x}{\sqrt{x^{2}+a^{2}}}&=\int\frac{a\sec^{2}t}{a\sec t}\mathrm{d}t=\int\sec t\,\mathrm{d}t=\int\frac{\mathrm{d}t}{\cos t}\\
&=\ln|\sec t+\tan t|+C_{1}\\
&=\ln\left|\frac{x}{a}+\frac{\sqrt{x^{2}+a^{2}}}{a}\right|+C_{1}\\
&=\ln(x+\sqrt{x^{2}+a^{2}})+C.
\end{align*}
\end{solution}

\begin{solution}[解法 3]
\begin{align*}
\int\frac{\mathrm{d}x}{\sqrt{x^{2}+a^{2}}}
&\xlongequal{x=a\sinh t}\int\frac{a\cosh t}{a\cosh t}\mathrm{d}t=\int\mathrm{d}t\\
&=t+C_{1}=\operatorname{arcsinh}\frac{x}{a}+C_{1}\\
&=\ln\left(\frac{x}{a}+\sqrt{\left(\frac{x}{a}\right)^{2}+1}\right)+C_{1}\\
&=\ln\left(x+\sqrt{x^{2}+a^{2}}\right)+C.
\end{align*}
其中 $\frac{x}{a}=\sinh t=\frac{\mathrm{e}^{t}-\mathrm{e}^{-t}}{2}$, $(\mathrm{e}^{t})^{2}-2\frac{x}{a}\mathrm{e}^{t}-1=0$. 由于 $\mathrm{e}^{t}>0$, 故
\begin{gather*}
\mathrm{e}^{t}=\frac{x}{a}+\sqrt{\left(\frac{x}{a}\right)^{2}+1},\\
\operatorname{arcsinh}\frac{x}{a}=t=\ln\left[\frac{x}{a}+\sqrt{\left(\frac{x}{a}\right)^{2}+1}\right].
\end{gather*}
\end{solution}

\begin{solution}[解法 4]
做变量代换 $\sqrt{x^{2}+a^{2}}=t-x$, 即 $x=\frac{t^{2}-a^{2}}{2t}$, 则
\begin{align*}
\int\frac{\mathrm{d}x}{\sqrt{x^{2}+a^{2}}}
&=\int\frac{1}{t-\frac{t^{2}-a^{2}}{2t}}\cdot\frac{2t^{2}-(t^{2}-a^{2})}{2t^{2}}\mathrm{d}t=\int\frac{\mathrm{d}t}{t}\\
&=\ln|t|+C=\ln(x+\sqrt{x^{2}+a^{2}})+C.
\end{align*}
\end{solution}

\begin{figure}[htbp]
\centering
\begin{tikzpicture}[x=1.1cm,y=1.1cm,>=Stealth]
  \draw[thick] (0,0) -- (2.4,0) node[midway,below] {$a$};
  \draw[thick] (0,0) -- (2.4,2.0);
  \draw[thick] (2.4,0) -- (2.4,2.0) node[midway,right] {$\sqrt{x^{2}+a^{2}}$};
  \draw (2.15,0) -- (2.15,0.25) -- (2.4,0.25);
  \draw (0.5,0) arc (0:39.8:0.5);
  \node at (0.62,0.18) {$t$};
  \node[above] at (0.85,1.35) {$x$};
  \node[right] at (3.0,1.7) {$\tan t=\frac{x}{a}$};
  \node[right] at (3.0,1.1) {$\sec t=\frac{\sqrt{x^{2}+a^{2}}}{a}$};
\end{tikzpicture}
\caption{}\label{fig:5.2.2}
\end{figure}

\begin{example}\label{ex:5.2.5}
求 $\int\sqrt{a^{2}-x^{2}}\,\mathrm{d}x$.
\end{example}

\begin{solution}
如图~\ref{fig:5.2.3} 所示.
\begin{align*}
\int\sqrt{a^{2}-x^{2}}\,\mathrm{d}x
&\xlongequal{x=a\sin t}\int a\cos t\,\mathrm{d}(a\sin t)=a^{2}\int\cos^{2}t\,\mathrm{d}t\\
&=\frac{a^{2}}{2}\int(1+\cos 2t)\mathrm{d}t=\frac{a^{2}}{2}\left(t+\frac{1}{2}\sin 2t\right)+C\\
&=\frac{a^{2}}{2}(t+\sin t\cos t)+C\\
&=\frac{a^{2}}{2}\left(\arcsin\frac{x}{a}+\frac{x}{a}\cdot\frac{\sqrt{a^{2}-x^{2}}}{a}\right)+C\\
&=\frac{1}{2}\left(a^{2}\arcsin\frac{x}{a}+x\sqrt{a^{2}-x^{2}}\right)+C.
\end{align*}
\end{solution}

\begin{figure}[htbp]
\centering
\begin{tikzpicture}[x=1.1cm,y=1.1cm,>=Stealth]
  \draw[thick] (0,0) -- (2.6,0) node[midway,below] {$\sqrt{a^{2}-x^{2}}$};
  \draw[thick] (0,0) -- (2.6,2.2);
  \draw[thick] (2.6,0) -- (2.6,2.2) node[midway,right] {$x$};
  \draw (2.35,0) -- (2.35,0.25) -- (2.6,0.25);
  \draw (0.55,0) arc (0:40.2:0.55);
  \node at (0.68,0.2) {$t$};
  \node[above] at (1.3,1.15) {$a$};
\end{tikzpicture}
\caption{}\label{fig:5.2.3}
\end{figure}

\begin{example}\label{ex:5.2.6}
求 $\int\frac{\mathrm{d}x}{(x^{2}+a^{2})^{3/2}}$ ($a>0$).
\end{example}

\begin{solution}[解法 1]
如图~\ref{fig:5.2.4} 所示.
\begin{align*}
\int\frac{\mathrm{d}x}{(x^{2}+a^{2})^{3/2}}
&\xlongequal{x=a\tan t}\int\frac{a\sec^{2}t}{a^{3}\sec^{3}t}\mathrm{d}t
=\frac{1}{a^{2}}\int\cos t\,\mathrm{d}t=\frac{1}{a^{2}}\sin t+C\\
&=\frac{x}{a^{2}\sqrt{x^{2}+a^{2}}}+C.
\end{align*}
\end{solution}

\begin{solution}[解法 2]
由分部积分法得到
\begin{align*}
\int\frac{\mathrm{d}x}{\sqrt{x^{2}+a^{2}}}
&=\frac{x}{\sqrt{x^{2}+a^{2}}}-\int x\left(-\frac{1}{2}\right)(x^{2}+a^{2})^{-\frac{3}{2}}\cdot 2x\,\mathrm{d}x\\
&=\frac{x}{\sqrt{x^{2}+a^{2}}}+\int\frac{(x^{2}+a^{2})-a^{2}}{(x^{2}+a^{2})^{3/2}}\mathrm{d}x\\
&=\frac{x}{\sqrt{x^{2}+a^{2}}}+\int\frac{\mathrm{d}x}{\sqrt{x^{2}+a^{2}}}-a^{2}\int\frac{\mathrm{d}x}{(x^{2}+a^{2})^{3/2}},
\end{align*}
移项化简得
\[
\int\frac{\mathrm{d}x}{(x^{2}+a^{2})^{3/2}}=\frac{x}{a^{2}\sqrt{x^{2}+a^{2}}}+C.
\]
\end{solution}

\begin{figure}[htbp]
\centering
\begin{tikzpicture}[x=1.1cm,y=1.1cm,>=Stealth]
  \draw[thick] (0,0) -- (2.4,0) node[midway,below] {$a$};
  \draw[thick] (0,0) -- (2.4,2.0);
  \draw[thick] (2.4,0) -- (2.4,2.0) node[midway,right] {$x$};
  \draw (2.15,0) -- (2.15,0.25) -- (2.4,0.25);
  \draw (0.5,0) arc (0:39.8:0.5);
  \node at (0.62,0.18) {$t$};
  \node[above] at (0.85,1.35) {$\sqrt{x^{2}+a^{2}}$};
\end{tikzpicture}
\caption{}\label{fig:5.2.4}
\end{figure}

\begin{example}\label{ex:5.2.7}
求 $\int\frac{\mathrm{d}x}{(x^{2}+a^{2})^{2}}$ ($a>0$).
\end{example}

\begin{solution}[解法 1]
如图~\ref{fig:5.2.5} 所示.
\begin{align*}
\int\frac{\mathrm{d}x}{(x^{2}+a^{2})^{2}}
&\xrightarrow{x=a\tan t}\int\frac{a\sec^{2}t}{a^{4}\sec^{4}t}\mathrm{d}t=\frac{1}{a^{3}}\int\cos^{2}t\,\mathrm{d}t\\
&=\frac{1}{2a^{3}}\int(1+\cos 2t)\mathrm{d}t=\frac{1}{2a^{3}}(t+\sin t\cos t)+C\\
&=\frac{1}{2a^{3}}\left(\arctan\frac{x}{a}+\frac{x}{\sqrt{x^{2}+a^{2}}}\cdot\frac{a}{\sqrt{x^{2}+a^{2}}}\right)+C\\
&=\frac{1}{2a^{3}}\left(\arctan\frac{x}{a}+\frac{ax}{x^{2}+a^{2}}\right)+C.
\end{align*}
\end{solution}

\begin{solution}[解法 2]
由
\begin{align*}
\int\frac{\mathrm{d}x}{x^{2}+a^{2}}&\xlongequal{\text{分部积分}}\frac{x}{x^{2}+a^{2}}-\int\frac{-2x}{(x^{2}+a^{2})^{2}}\mathrm{d}x\\
&=\frac{x}{x^{2}+a^{2}}+2\int\frac{(x^{2}+a^{2})-a^{2}}{(x^{2}+a^{2})^{2}}\mathrm{d}x\\
&=\frac{x}{x^{2}+a^{2}}+2\int\frac{\mathrm{d}x}{x^{2}+a^{2}}-2a^{2}\int\frac{\mathrm{d}x}{(x^{2}+a^{2})^{2}},
\end{align*}
得到
\begin{align*}
\int\frac{\mathrm{d}x}{(x^{2}+a^{2})^{2}}&=\frac{1}{2a^{2}}\cdot\frac{x}{x^{2}+a^{2}}+\frac{1}{2a^{2}}\int\frac{\mathrm{d}x}{x^{2}+a^{2}}\\
&=\frac{x}{2a^{2}(x^{2}+a^{2})}+\frac{1}{2a^{3}}\arctan\frac{x}{a}+C.
\end{align*}
\end{solution}

\begin{figure}[htbp]
\centering
\begin{tikzpicture}[x=1.1cm,y=1.1cm,>=Stealth]
  \draw[thick] (0,0) -- (2.4,0) node[midway,below] {$a$};
  \draw[thick] (0,0) -- (2.4,2.0);
  \draw[thick] (2.4,0) -- (2.4,2.0) node[midway,right] {$x$};
  \draw (2.15,0) -- (2.15,0.25) -- (2.4,0.25);
  \draw (0.5,0) arc (0:39.8:0.5);
  \node at (0.62,0.18) {$t$};
  \node[above] at (0.95,1.45) {$\sqrt{x^{2}+a^{2}}$};
\end{tikzpicture}
\caption{}\label{fig:5.2.5}
\end{figure}

\begin{example}\label{ex:5.2.8}
求 $\int\frac{\mathrm{d}x}{x^{2}\sqrt{x^{2}-1}}$.
\end{example}

\begin{solution}[解法 1]
如图~\ref{fig:5.2.6} 所示.
\begin{align*}
\int\frac{\mathrm{d}x}{x^{2}\sqrt{x^{2}-1}}
&\xlongequal{x=\sec t}\int\frac{\sec t\tan t}{\sec^{2}t\tan t}\mathrm{d}t\\
&=\int\cos t\,\mathrm{d}t=\sin t+C\\
&=\frac{\sqrt{x^{2}-1}}{x}+C.
\end{align*}
\end{solution}

\begin{solution}[解法 2]
\begin{align*}
\int\frac{\mathrm{d}x}{x^{2}\sqrt{x^{2}-1}}&=\int\frac{\mathrm{d}x}{x^{3}\sqrt{1-\frac{1}{x^{2}}}}
=\int\frac{1}{x}\cdot\frac{-1}{\sqrt{1-\frac{1}{x^{2}}}}\mathrm{d}\left(\frac{1}{x}\right)\\
&\xlongequal{u=\frac{1}{x}}\int\frac{-u}{\sqrt{1-u^{2}}}\mathrm{d}u=\sqrt{1-u^{2}}+C\\
&=\sqrt{1-\left(\frac{1}{x}\right)^{2}}+C=\frac{\sqrt{x^{2}-1}}{x}+C.
\end{align*}
或
\begin{align*}
\int\frac{\mathrm{d}x}{x^{2}\sqrt{x^{2}-1}}&=\int\frac{1}{x^{3}}\cdot\frac{\mathrm{d}x}{\sqrt{1-\frac{1}{x^{2}}}}
=\frac{1}{2}\int\frac{1}{\sqrt{1-\frac{1}{x^{2}}}}\mathrm{d}\left(1-\frac{1}{x^{2}}\right)\\
&\xlongequal{u=1-\frac{1}{x^{2}}}\int\frac{\mathrm{d}u}{2\sqrt{u}}=\sqrt{u}+C\\
&=\sqrt{1-\frac{1}{x^{2}}}+C=\frac{\sqrt{x^{2}-1}}{x}+C.
\end{align*}
\end{solution}

\begin{figure}[htbp]
\centering
\begin{tikzpicture}[x=1.4cm,y=1.4cm,>=Stealth]
  \draw[thick] (0,0) -- (2.0,0) node[midway,below] {$1$};
  \draw[thick] (0,0) -- (2.0,1.7);
  \draw[thick] (2.0,0) -- (2.0,1.7) node[midway,right] {$\sqrt{x^{2}-1}$};
  \draw (1.75,0) -- (1.75,0.25) -- (2.0,0.25);
  \draw (0.5,0) arc (0:40.4:0.5);
  \node at (0.62,0.18) {$t$};
  \node[above] at (0.75,1.25) {$x$};
\end{tikzpicture}
\caption{}\label{fig:5.2.6}
\end{figure}

\begin{example}\label{ex:5.2.9}
求 $\int\frac{\mathrm{d}u}{\sqrt{u}+\sqrt[3]{u}}$.
\end{example}

\begin{solution}
\begin{align*}
\int\frac{\mathrm{d}u}{\sqrt{u}+\sqrt[3]{u}}&\xlongequal{u=x^{6}}\int\frac{6x^{5}}{x^{3}+x^{2}}\mathrm{d}x
=6\int\frac{x^{3}}{x+1}\mathrm{d}x\\
&=6\int\left(x^{2}-x+1-\frac{1}{x+1}\right)\mathrm{d}x\\
&=6\left(\frac{x^{3}}{3}-\frac{x^{2}}{2}+x-\ln|x+1|\right)+C\\
&=2\sqrt{u}-3\sqrt[3]{u}+6\sqrt[6]{u}-6\ln(\sqrt[6]{u}+1)+C.
\end{align*}
\end{solution}

除了换元 (变量代换) 积分法外, 还有一种重要的分部积分法, 它有 4 种不同的类型: 升幂、降幂、循环、递推.

\begin{example}\label{ex:5.2.10}
利用升幂的方法求下列不定积分:
\begin{enumerate}[label=(\arabic*)]
  \item $\int x^{2}\arctan x\,\mathrm{d}x$;
  \item $\int\arctan x\,\mathrm{d}x$;
  \item $\int\frac{\ln x}{x^{2}}\mathrm{d}x$;
  \item $\int\ln x\,\mathrm{d}x$;
  \item $\int x\ln^{2}x\,\mathrm{d}x$;
  \item $\int x^{3}\ln x\,\mathrm{d}x$.
\end{enumerate}
\end{example}

\begin{solution}
(1)
\begin{align*}
\int x^{2}\arctan x\,\mathrm{d}x&=\frac{1}{3}\int\arctan x\,\mathrm{d}x^{3}
=\frac{1}{3}\left(x^{3}\arctan x-\int\frac{x^{3}}{1+x^{2}}\mathrm{d}x\right)\\
&=\frac{1}{3}\left[x^{3}\arctan x-\int\left(x-\frac{x}{1+x^{2}}\right)\mathrm{d}x\right]\\
&=\frac{1}{3}\left[x^{3}\arctan x-\int x\,\mathrm{d}x+\frac{1}{2}\int\frac{\mathrm{d}(1+x^{2})}{1+x^{2}}\right]\\
&=\frac{1}{3}x^{3}\arctan x-\frac{x^{2}}{6}+\frac{1}{6}\ln(1+x^{2})+C,
\end{align*}
其中 $x^{2}\arctan x$ 中的 $x^{2}$ 升幂为 $\mathrm{d}x^{3}$ 中的 $x^{3}$. 通过分部积分将一个不易算的积分\linebreak[4]$\int x^{2}\arctan x\,\mathrm{d}x$ 化为一个易算的积分 $\int\frac{x^{3}}{1+x^{2}}\mathrm{d}x$.

(2)
\begin{align*}
\int\arctan x\,\mathrm{d}x&=x\arctan x-\int\frac{x}{1+x^{2}}\mathrm{d}x=x\arctan x-\frac{1}{2}\int\frac{\mathrm{d}(1+x^{2})}{1+x^{2}}\\
&=x\arctan x-\frac{1}{2}\ln(1+x^{2})+C.
\end{align*}
其中 $\arctan x=x^{0}\arctan x$ 中的 $x^{0}$ 已经升幂为 $\mathrm{d}x$ 中的 $x=x^{1}$.

(3)
\begin{align*}
\int\frac{\ln x}{x^{2}}\mathrm{d}x&=-\int\ln x\,\mathrm{d}\left(\frac{1}{x}\right)=-\frac{\ln x}{x}+\int\frac{1}{x}\cdot\frac{1}{x}\mathrm{d}x\\
&=-\frac{\ln x}{x}-\frac{1}{x}+C.
\end{align*}
(4)
\begin{align*}
\int\ln x\,\mathrm{d}x&=x\ln x-\int x\cdot\frac{1}{x}\mathrm{d}x=x\ln x-\int\mathrm{d}x\\
&=x\ln x-x+C.
\end{align*}
(5)
\begin{align*}
\int x\ln^{2}x\,\mathrm{d}x&=\frac{1}{2}\int\ln^{2}x\,\mathrm{d}x^{2}=\frac{1}{2}\left(x^{2}\ln^{2}x-\int x^{2}\cdot 2\ln x\cdot\frac{1}{x}\mathrm{d}x\right)\\
&=\frac{x^{2}}{2}\ln^{2}x-\int x\ln x\,\mathrm{d}x=\frac{x^{2}}{2}\ln^{2}x-\frac{1}{2}\int\ln x\,\mathrm{d}x^{2}\\
&=\frac{x^{2}}{2}\ln^{2}x-\frac{1}{2}\left(x^{2}\ln x-\int x^{2}\cdot\frac{1}{x}\mathrm{d}x\right)\\
&=\frac{x^{2}}{2}\ln^{2}x-\frac{x^{2}}{2}\ln x+\frac{x^{2}}{4}+C.
\end{align*}
(6)
\begin{align*}
\int x^{3}\ln x\,\mathrm{d}x&=\int\ln x\,\mathrm{d}\frac{x^{4}}{4}=\frac{x^{4}}{4}\ln x-\int\frac{x^{4}}{4}\cdot\frac{1}{x}\mathrm{d}x\\
&=\frac{x^{4}}{4}\ln x-\frac{1}{4}\int x^{3}\mathrm{d}x\\
&=\frac{x^{4}}{4}\ln x-\frac{x^{4}}{16}+C.
\end{align*}
\end{solution}

\begin{example}\label{ex:5.2.11}
利用降幂的方法求下列不定积分:
\begin{enumerate}[label=(\arabic*)]
  \item $\int x\cos x\,\mathrm{d}x$;
  \item $\int x^{2}\mathrm{e}^{-x}\mathrm{d}x$;
  \item $\int x^{2}\sin x\,\mathrm{d}x$.
\end{enumerate}
\end{example}

\begin{solution}
(1)
\begin{align*}
\int x\cos x\,\mathrm{d}x&=\int x\,\mathrm{d}\sin x=x\sin x-\int\sin x\,\mathrm{d}x\\
&=x\sin x+\cos x+C,
\end{align*}
其中 $x\cos x$ 中的 $x$ 降幂为 $\sin x=x^{0}\sin x$ 中的 $x^{0}$. 通过分部积分将一个不易算的积分 $\int x\cos x\,\mathrm{d}x$ 化为一个易算的积分 $\int\sin x\,\mathrm{d}x$.

(2)
\begin{align*}
\int x^{2}\mathrm{e}^{-x}\mathrm{d}x&=-\int x^{2}\mathrm{d}\mathrm{e}^{-x}=-x^{2}\mathrm{e}^{-x}+\int\mathrm{e}^{-x}\cdot 2x\,\mathrm{d}x\\
&=-x^{2}\mathrm{e}^{-x}-2\int x\,\mathrm{d}\mathrm{e}^{-x}\\
&=-x^{2}\mathrm{e}^{-x}-2\left(x\mathrm{e}^{-x}-\int\mathrm{e}^{-x}\mathrm{d}x\right)\\
&=-x^{2}\mathrm{e}^{-x}-2x\mathrm{e}^{-x}-2\mathrm{e}^{-x}+C\\
&=-\mathrm{e}^{-x}(x^{2}+x+2)+C.
\end{align*}
(3)
\begin{align*}
\int x^{2}\sin x\,\mathrm{d}x&=-\int x^{2}\mathrm{d}\cos x=-x^{2}\cos x+\int 2x\cos x\,\mathrm{d}x\\
&=-x^{2}\cos x+2\int x\,\mathrm{d}\sin x\\
&=-x^{2}\cos x+2\left(x\sin x-\int\sin x\,\mathrm{d}x\right)\\
&=-x^{2}\cos x+2x\sin x+2\cos x+C.
\end{align*}
\end{solution}

\begin{example}\label{ex:5.2.12}
利用循环的方法求 $I=\int\mathrm{e}^{ax}\cos bx\,\mathrm{d}x$, $J=\int\mathrm{e}^{ax}\sin bx\,\mathrm{d}x$.
\end{example}

\begin{solution}[解法 1]
当 $a\neq 0$ 时,
\begin{align*}
I&=\int\mathrm{e}^{ax}\cos bx\,\mathrm{d}x=\frac{1}{a}\int\cos bx\,\mathrm{d}\mathrm{e}^{ax}
=\frac{1}{a}\left(\mathrm{e}^{ax}\cos bx+b\int\mathrm{e}^{ax}\sin bx\,\mathrm{d}x\right)\\
&=\frac{1}{a}\mathrm{e}^{ax}\cos bx+\frac{b}{a}J,\\
J&=\int\mathrm{e}^{ax}\sin bx\,\mathrm{d}x=\frac{1}{a}\int\sin bx\,\mathrm{d}\mathrm{e}^{ax}
=\frac{1}{a}\left(\mathrm{e}^{ax}\sin bx-b\int\mathrm{e}^{ax}\cos bx\,\mathrm{d}x\right)\\
&=\frac{1}{a}\mathrm{e}^{ax}\sin bx-\frac{b}{a}I.
\end{align*}
由此得到
\[
\begin{cases}
aI-bJ=\mathrm{e}^{ax}\cos bx,\\
bI+aJ=\mathrm{e}^{ax}\sin bx.
\end{cases}
\]
解此方程得
\begin{gather*}
I=\int\mathrm{e}^{ax}\cos bx\,\mathrm{d}x=\frac{b\sin bx+a\cos bx}{a^{2}+b^{2}}\mathrm{e}^{ax}+C_{1},\\
J=\int\mathrm{e}^{ax}\sin bx\,\mathrm{d}x=\frac{a\sin bx-b\cos bx}{a^{2}+b^{2}}\mathrm{e}^{ax}+C_{2}.
\end{gather*}
注意 $a=0$, $b\neq 0$ 时, 上述公式仍成立. 当 $a=b=0$ 时,
\[
I=\int\mathrm{d}x=x+C_{1},\quad J=\int 0\,\mathrm{d}x=C_{2}.
\]
\end{solution}

\begin{solution}[解法 2]
当 $a\neq 0$ 时,
\begin{align*}
I&=\int\mathrm{e}^{ax}\cos bx\,\mathrm{d}x=\frac{1}{a}\int\cos bx\,\mathrm{d}\mathrm{e}^{ax}\\
&=\frac{1}{a}\left(\mathrm{e}^{ax}\cos bx+b\int\mathrm{e}^{ax}\sin bx\,\mathrm{d}x\right)\\
&=\frac{1}{a}\left(\mathrm{e}^{ax}\cos bx+\frac{b}{a}\int\sin bx\,\mathrm{d}\mathrm{e}^{ax}\right)\\
&=\frac{1}{a}\mathrm{e}^{ax}\cos bx+\frac{b}{a^{2}}\left(\mathrm{e}^{ax}\sin bx-b\int\mathrm{e}^{ax}\cos bx\,\mathrm{d}x\right)\\
&=\frac{b\sin bx+a\cos bx}{a^{2}}\mathrm{e}^{ax}-\frac{b^{2}}{a^{2}}\int\mathrm{e}^{ax}\cos bx\,\mathrm{d}x\\
&=\frac{b\sin bx+a\cos bx}{a^{2}}\mathrm{e}^{ax}-\frac{b^{2}}{a^{2}}I,
\end{align*}
由于右边出现 $I$ (循环), 移项化简得
\begin{align*}
I&=\int\mathrm{e}^{ax}\cos bx\,\mathrm{d}x=\frac{a^{2}}{a^{2}+b^{2}}\cdot\frac{b\sin bx+a\cos bx}{a^{2}}\mathrm{e}^{ax}+C_{1}\\
&=\frac{b\sin bx+a\cos bx}{a^{2}+b^{2}}\mathrm{e}^{ax}+C_{1}.
\end{align*}
同理可得
\[
J=\frac{a\sin bx-b\cos bx}{a^{2}+b^{2}}\mathrm{e}^{ax}+C_{2}.
\]
或者, 当 $b\neq 0$ 时,
\begin{align*}
J&=\int\mathrm{e}^{ax}\sin bx\,\mathrm{d}x=\frac{1}{b}(aI-\mathrm{e}^{ax}\cos bx)\\
&=\frac{1}{b}\cdot\frac{ab\sin bx+a^{2}\cos bx-a^{2}\cos bx-b^{2}\cos bx}{a^{2}+b^{2}}\mathrm{e}^{ax}+\frac{a}{b}C_{1}\\
&=\frac{a\sin bx-b\cos bx}{a^{2}+b^{2}}\mathrm{e}^{ax}+C_{2}.
\end{align*}
\end{solution}

\begin{example}\label{ex:5.2.13}
求 $\int\sqrt{x^{2}+a^{2}}\,\mathrm{d}x$ ($a>0$).
\end{example}

\begin{solution}[解法 1]
应用分部积分得到
\begin{align*}
\int\sqrt{x^{2}+a^{2}}\,\mathrm{d}x&=x\sqrt{x^{2}+a^{2}}-\int x\cdot\frac{2x}{2\sqrt{x^{2}+a^{2}}}\mathrm{d}x\\
&=x\sqrt{x^{2}+a^{2}}-\int\frac{(x^{2}+a^{2})-a^{2}}{\sqrt{x^{2}+a^{2}}}\mathrm{d}x\\
&=x\sqrt{x^{2}+a^{2}}-\int\sqrt{x^{2}+a^{2}}\,\mathrm{d}x+a^{2}\int\frac{\mathrm{d}x}{\sqrt{x^{2}+a^{2}}},
\end{align*}
这就出现了循环, 由此解得
\begin{align*}
\int\sqrt{x^{2}+a^{2}}\,\mathrm{d}x&=\frac{x}{2}\sqrt{x^{2}+a^{2}}+\frac{a^{2}}{2}\int\frac{\mathrm{d}x}{\sqrt{x^{2}+a^{2}}}\\
&=\frac{x}{2}\sqrt{x^{2}+a^{2}}+\frac{a^{2}}{2}\ln(x+\sqrt{x^{2}+a^{2}})+C.
\end{align*}
\end{solution}

\begin{solution}[解法 2]
\begin{align*}
\int\sqrt{x^{2}+a^{2}}\,\mathrm{d}x&\xlongequal{x=a\sinh u}\int\sqrt{a^{2}\sinh^{2}u+a^{2}}\,a\cosh u\,\mathrm{d}u\\
&=\int a^{2}\cosh^{2}u\,\mathrm{d}u=a^{2}\int\frac{1+\cosh 2u}{2}\mathrm{d}u\\
&=\frac{a^{2}}{2}\left(u+\frac{\sinh 2u}{2}\right)+C_{1}\\
&=\frac{a^{2}}{2}(\sinh u\cosh u+u)+C_{1}\\
&=\frac{x}{2}\sqrt{x^{2}+a^{2}}+\frac{a^{2}}{2}\ln\left(x+\sqrt{x^{2}+a^{2}}\right)+C.
\end{align*}
\end{solution}

\begin{solution}[解法 3]
\begin{align*}
\int\sqrt{x^{2}+a^{2}}\,\mathrm{d}x&\xlongequal{x=a\tan t}\int\sqrt{a^{2}\tan^{2}t+a^{2}}\,a\sec^{2}t\,\mathrm{d}t\\
&=a^{2}\int\frac{\mathrm{d}t}{\cos^{3}t}=a^{2}\int\frac{\cos t}{(1-\sin^{2}t)^{2}}\mathrm{d}t
=a^{2}\int\frac{\mathrm{d}\sin t}{(1-\sin^{2}t)^{2}}\\
&\xlongequal{u=\sin t}a^{2}\int\frac{\mathrm{d}u}{(1-u^{2})^{2}}=\frac{a^{2}}{4}\int\left(\frac{1}{1-u}+\frac{1}{1+u}\right)^{2}\mathrm{d}u\\
&=\frac{a^{2}}{4}\int\left[\frac{1}{(1-u)^{2}}+\frac{2}{(1-u)(1+u)}+\frac{1}{(1+u)^{2}}\right]\mathrm{d}u\\
&=\frac{a^{2}}{4}\int\left[\frac{1}{(1-u)^{2}}+\frac{1}{1-u}+\frac{1}{1+u}+\frac{1}{(1+u)^{2}}\right]\mathrm{d}u\\
&=\frac{a^{2}}{4}\left(\frac{1}{1-u}+\ln\left|\frac{1+u}{1-u}\right|-\frac{1}{1+u}\right)+C_{1}\\
&=\frac{a^{2}}{4}\left(\frac{2u}{1-u^{2}}+\ln\left|\frac{1+u}{1-u}\right|\right)+C_{1}\\
&=\frac{a^{2}}{4}\left(\frac{2\sin t}{1-\sin^{2}t}+\ln\frac{1+\sin t}{1-\sin t}\right)+C_{1}\\
&=\frac{1}{2}a\tan t\cdot\frac{a}{\cos t}+\frac{a^{2}}{4}\ln\frac{1+\frac{x}{\sqrt{x^{2}+a^{2}}}}{1-\frac{x}{\sqrt{x^{2}+a^{2}}}}+C_{1}\\
&=\frac{x}{2}\sqrt{x^{2}+a^{2}}+\frac{a^{2}}{4}\ln\frac{\sqrt{x^{2}+a^{2}}+x}{\sqrt{x^{2}+a^{2}}-x}+C_{1}\\
&=\frac{x}{2}\sqrt{x^{2}+a^{2}}+\frac{a^{2}}{4}\ln\frac{(x+\sqrt{x^{2}+a^{2}})^{2}}{a^{2}}+C_{1}\\
&=\frac{x}{2}\sqrt{x^{2}+a^{2}}+\frac{a^{2}}{2}\ln(x+\sqrt{x^{2}+a^{2}})+C.
\end{align*}
\end{solution}

\begin{example}\label{ex:5.2.14}
利用递推的方法求下面积分.
\[
I_{n}=\int\sin^{n}x\,\mathrm{d}x,\quad J_{n}=\int\cos^{n}x\,\mathrm{d}x,\quad n\in\mathbb{N}.
\]
\end{example}

\begin{solution}
\begin{align*}
I_{n}&=\int\sin^{n}x\,\mathrm{d}x=-\int\sin^{n-1}x\,\mathrm{d}\cos x\\
&\xlongequal{\text{分部积分}}-\sin^{n-1}x\cos x+\int(n-1)\sin^{n-2}x\cos^{2}x\,\mathrm{d}x\\
&=-\sin^{n-1}x\cos x+(n-1)\int\sin^{n-2}x(1-\sin^{2}x)\mathrm{d}x\\
&=-\sin^{n-1}x\cos x+(n-1)I_{n-2}-(n-1)I_{n}.
\end{align*}
由此得到递推公式
\[
I_{n}=\int\sin^{n}x\,\mathrm{d}x=-\frac{1}{n}\sin^{n-1}x\cos x+\frac{n-1}{n}I_{n-2}.
\]
\begin{align*}
J_{n}&=\int\cos^{n}x\,\mathrm{d}x=\int\cos^{n-1}x\,\mathrm{d}\sin x\\
&\xlongequal{\text{分部积分}}\cos^{n-1}x\sin x+(n-1)\int\cos^{n-2}x(1-\cos^{2}x)\mathrm{d}x\\
&=\cos^{n-1}x\sin x+(n-1)J_{n-2}-(n-1)J_{n}.
\end{align*}
同样, 由此得到递推公式
\[
J_{n}=\int\cos^{n}x\,\mathrm{d}x=\frac{1}{n}\cos^{n-1}x\sin x+\frac{n-1}{n}J_{n-2}.
\]
从第一个递推公式得到,
\begin{align*}
I_{4}&=\int\sin^{4}x\,\mathrm{d}x=-\frac{1}{4}\sin^{3}x\cos x+\frac{3}{4}I_{2}\\
&=-\frac{1}{4}\sin^{3}x\cos x+\frac{3}{4}\left(-\frac{1}{2}\sin x\cos x+\frac{1}{2}I_{0}\right)\\
&=-\frac{1}{4}\sin^{3}x\cos x-\frac{3}{8}\sin x\cos x+\frac{3}{8}x+C.
\end{align*}
\end{solution}

\begin{exblock}{练习题 5.2}
  \item 应用换元积分法求下列不定积分:
  \begin{enumerate}[label=(\arabic*)]
    \item $\int\frac{\mathrm{d}x}{(a+bx)^{2}}$ ($a,b$ 为常数, $ab\neq 0$);
    \item $\int\frac{x^{2}+2}{(x+1)^{3}}\mathrm{d}x$;
    \item $\int\frac{\mathrm{d}x}{2x+1}$;
    \item $\int(1+x)^{n}\mathrm{d}x$;
    \item $\int\left(\frac{1}{\sqrt{3-x^{2}}}+\frac{1}{\sqrt{1-3x^{2}}}\right)\mathrm{d}x$;
    \item $\int 2^{2x+3}\mathrm{d}x$;
    \item $\int\frac{\mathrm{d}x}{\sqrt{x}(1+x)}$;
    \item $\int\frac{\mathrm{d}x}{\sqrt[3]{7-5x}}$;
    \item $\int\frac{x}{\sqrt{1-x^{2}}}\mathrm{d}x$;
    \item $\int\frac{x}{4+x^{2}}\mathrm{d}x$;
    \item $\int\frac{x^{4}}{(1-x^{5})^{3}}\mathrm{d}x$;
    \item $\int\frac{x^{3}}{x^{8}-2}\mathrm{d}x$;
    \item $\int\frac{\mathrm{d}x}{x(1+x)}$;
    \item $\int\frac{2x-3}{x^{2}-3x+8}\mathrm{d}x$;
    \item $\int\cos(3x+4)\mathrm{d}x$;
    \item $\int\frac{\mathrm{d}x}{\sin^{2}\left(2x+\frac{\pi}{4}\right)}$;
    \item $\int\frac{\mathrm{d}x}{1+\cos x}$;
    \item $\int\frac{\mathrm{d}x}{1+\sin x}$;
    \item $\int\csc x\,\mathrm{d}x$;
    \item $\int\cot x\,\mathrm{d}x$;
    \item $\int\cos^{5}x\,\mathrm{d}x$;
    \item $\int\frac{\mathrm{d}x}{\sin x}$;
    \item $\int\frac{\mathrm{d}x}{\sin x\cos x}$;
    \item $\int\frac{1}{x^{2}}\sin\frac{1}{x}\mathrm{d}x$;
    \item $\int\frac{(\arctan x)^{2}}{1+x^{2}}\mathrm{d}x$;
    \item $\int\cos ax\sin bx\,\mathrm{d}x$;
    \item $\int\cos ax\cos bx\,\mathrm{d}x$;
    \item $\int\sin ax\sin bx\,\mathrm{d}x$;
    \item $\int\sin^{5}x\,\mathrm{d}x$;
    \item $\int\sin^{4}x\,\mathrm{d}x$;
    \item $\int\frac{\mathrm{d}x}{\cos x+\sin x}$;
    \item $\int\frac{\mathrm{d}x}{a\cos x+b\sin x}$;
    \item $\int\frac{\cos x}{\sqrt{2+\cos 2x}}\mathrm{d}x$;
    \item $\int\frac{\cos x\sin x}{1+\sin^{4}x}\mathrm{d}x$;
    \item $\int\frac{\mathrm{d}x}{1+\varepsilon\cos x}$ ($\varepsilon\in(0,1)$);
    \item $\int\frac{\mathrm{d}x}{2-\sin^{2}x}$;
    \item $\int\frac{\cos x\sin x}{a^{2}\cos^{2}x+b^{2}\sin^{2}x}\mathrm{d}x$;
    \item $\int\frac{\cos x+\sin x}{\sqrt[3]{\sin x}-\cos x}\mathrm{d}x$;
    \item $\int\frac{\mathrm{d}x}{\cosh x}$;
    \item $\int\frac{\mathrm{d}x}{x\ln x(\ln\ln x)}$;
    \item $\int\frac{\ln^{2}x}{x}\mathrm{d}x$;
    \item $\int x\mathrm{e}^{-x^{2}}\mathrm{d}x$;
    \item $\int\frac{\ln x}{x\sqrt{1+\ln x}}\mathrm{d}x$;
    \item $\int\frac{x^{3}}{1-x^{2}}\mathrm{d}x$;
    \item $\int\mathrm{e}^{\sqrt{x}}\mathrm{d}x$;
    \item $\int x^{3}\mathrm{e}^{-x^{2}}\mathrm{d}x$;
    \item $\int\frac{x\,\mathrm{d}x}{\sqrt{1+x^{2}}+(1+x^{2})^{3/2}}$;
    \item $\int\frac{\mathrm{d}x}{x\sqrt{x^{2}+1}}$;
    \item $\int\frac{x^{2}}{\sqrt{x^{2}+a^{2}}}\mathrm{d}x$;
    \item $\int\arctan\sqrt{x}\,\mathrm{d}x$;
    \item $\int\frac{\mathrm{d}x}{\mathrm{e}^{x}+\mathrm{e}^{-x}}$;
    \item $\int\frac{x^{\frac{n}{2}}}{\sqrt{1+x^{n+2}}}\mathrm{d}x$;
    \item $\int\frac{\sqrt{x}}{1-\sqrt[3]{x}}\mathrm{d}x$;
    \item $\int\frac{\mathrm{d}x}{(x^{2}+a^{2})^{\frac{3}{2}}}$;
    \item $\int\frac{\sqrt{x+1}-1}{\sqrt{x+1}+1}\mathrm{d}x$;
    \item $\int\frac{\mathrm{d}x}{\cos^{4}x}$;
    \item $\int\sin^{4}x\,\mathrm{d}x$.
  \end{enumerate}

  \item 应用分部积分法求下列不定积分:
  \begin{enumerate}[label=(\arabic*)]
    \item $\int\arcsin x\,\mathrm{d}x$;
    \item $\int\ln x\,\mathrm{d}x$;
    \item $\int x^{2}\cos x\,\mathrm{d}x$;
    \item $\int\frac{\ln x}{x^{3}}\mathrm{d}x$;
    \item $\int\ln^{2}x\,\mathrm{d}x$;
    \item $\int x\arctan x\,\mathrm{d}x$;
    \item $\int(\arcsin x)^{2}\mathrm{d}x$;
    \item $\int\sec^{3}x\,\mathrm{d}x$;
    \item $\int x\arcsin x\,\mathrm{d}x$;
    \item $\int\left(\ln\ln x+\frac{1}{\ln x}\right)\mathrm{d}x$;
    \item $\int\arctan x\,\mathrm{d}x$;
    \item $\int x^{2}\arccos x\,\mathrm{d}x$;
    \item $\int x^{2}\cos x\,\mathrm{d}x$;
    \item $\int\frac{x}{\cos^{2}x}\mathrm{d}x$;
    \item $\int\frac{\arctan x}{x^{2}}\mathrm{d}x$;
    \item $\int x^{2}\cosh x\,\mathrm{d}x$;
    \item $\int\cos(\ln x)\mathrm{d}x$;
    \item $\int\sin(\ln x)\mathrm{d}x$;
    \item $\int\sqrt{x}\ln^{2}x\,\mathrm{d}x$;
    \item $\int\ln(x+\sqrt{1+x^{2}})\mathrm{d}x$.
  \end{enumerate}

  \item 根据公式 $\cosh^{2}t-\sinh^{2}t=1$, 用变量代换 $x=a\sinh t$ 或 $x=a\cosh t$ 求下列不定积分:
  \begin{enumerate}[label=(\arabic*)]
    \item $\int\frac{\mathrm{d}x}{\sqrt{x^{2}+a^{2}}}$;
    \item $\int\sqrt{x^{2}+a^{2}}\,\mathrm{d}x$;
    \item $\int(x^{2}+a^{2})^{-\frac{3}{2}}\mathrm{d}x$;
    \item $\int\frac{x^{2}}{\sqrt{x^{2}-a^{2}}}\mathrm{d}x$;
    \item $\int\sqrt{x^{2}-a^{2}}\,\mathrm{d}x$;
    \item $\int\frac{x\ln(x+\sqrt{1+x^{2}})}{\sqrt{1+x^{2}}}\mathrm{d}x$.
  \end{enumerate}

  \item 求函数 $f(x)$, 设
  \begin{enumerate}[label=(\arabic*)]
    \item $f'(x^{2})=\frac{1}{x}$ ($x>0$);
    \item $f'(\sin^{2}x)=\cos^{2}x$.
  \end{enumerate}

  \item 求不定积分:
  \begin{enumerate}[label=(\arabic*)]
    \item $\int\left(1-\frac{2}{x}\right)^{2}\mathrm{e}^{x}\mathrm{d}x$;
    \item $\int\frac{x\mathrm{e}^{x}}{(1+x)^{2}}\mathrm{d}x$;
    \item $\int\arctan(1+\sqrt{x})\mathrm{d}x$.
  \end{enumerate}

  \item 求下列不定积分:
  \begin{enumerate}[label=(\arabic*)]
    \item $\int[f(x)]^{\alpha}f'(x)\mathrm{d}x$ ($\alpha\neq -1$);
    \item $\int\frac{f'(x)}{1+[f(x)]^{2}}\mathrm{d}x$;
    \item $\int\frac{f'(x)}{f(x)}\mathrm{d}x$;
    \item $\int\mathrm{e}^{f(x)}f'(x)\mathrm{d}x$.
  \end{enumerate}

  \item 计算: (1) $\int xf''(x)\mathrm{d}x$; (2) $\int\mathrm{e}^{\sin x}\sin 2x\,\mathrm{d}x$.

  \item 设 $P$ 为 $n$ 次多项式, 试通过 $P$ 的各阶导数来表示不定积分 $\int P(x)\mathrm{e}^{ax}\mathrm{d}x$.

  \item 推导出不定积分 $\int\ln^{n}x\,\mathrm{d}x$ 的递推公式, $n\in\mathbb{N}$.

  \item 设 $I_{n}=\int\tan^{n}x\,\mathrm{d}x$, $n=2,3,\cdots$. 证明: 递推公式 $I_{n}=\frac{1}{n-1}\tan^{n-1}x-I_{n-2}$.

  \item 设 $I(m,n)=\int\cos^{m}x\sin^{n}x\,\mathrm{d}x$, 则当 $m+n\neq 0$ 时, 证明: 递推公式
  \begin{align*}
  I(m,n)&=\frac{\cos^{m-1}x\sin^{n+1}x}{m+n}+\frac{m-1}{m+n}I(m-2,n)\\
  &=-\frac{\cos^{m+1}x\sin^{n-1}x}{m+n}+\frac{n-1}{m+n}I(m,n-2),\quad n,m=2,3,\cdots.
\end{align*}

  \item 应用题 10 中的递推公式计算:
  \begin{enumerate}[label=(\arabic*)]
    \item $\int\tan^{3}x\,\mathrm{d}x$;
    \item $\int\tan^{4}x\,\mathrm{d}x$.
  \end{enumerate}

  \item 应用题 11 中的递推公式计算:
  \begin{enumerate}[label=(\arabic*)]
    \item $\int\cos^{2}x\sin^{4}x\,\mathrm{d}x$;
    \item $\int\cos^{4}x\sin^{2}x\,\mathrm{d}x$.
  \end{enumerate}

  \item 导出下列不定积分对于正整数 $n$ 的递推公式:
  \begin{enumerate}[label=(\arabic*)]
    \item $I_{n}=\int x^{n}\mathrm{e}^{kx}\,\mathrm{d}x$;
    \item $I_{n}=\int\ln^{n}x\,\mathrm{d}x$;
    \item $I_{n}=\int(\arcsin x)^{n}\mathrm{d}x$;
    \item $I_{n}=\int\mathrm{e}^{x}\sin^{n}x\,\mathrm{d}x$.
  \end{enumerate}

  \item 应用题 14 中的递推公式计算:
  \begin{enumerate}[label=(\arabic*)]
    \item $\int x^{3}\mathrm{e}^{2x}\mathrm{d}x$;
    \item $\int\ln^{3}x\,\mathrm{d}x$;
    \item $\int(\arcsin x)^{3}\mathrm{d}x$;
    \item $\int\mathrm{e}^{x}\sin^{3}x\,\mathrm{d}x$.
  \end{enumerate}
\end{exblock}

\section{有理函数的不定积分、可化为有理函数的不定积分}\label{sec:5.3}

由两个实多项式函数之商所表示的函数称为 (实) 有理函数, 它的一般形式为
\[
R(x)=\frac{P(x)}{Q(x)}
=\frac{\alpha_{0}x^{n}+\alpha_{1}x^{n-1}+\cdots+\alpha_{n-1}x+\alpha_{n}}
{\beta_{0}x^{m}+\beta_{1}x^{m-1}+\cdots+\beta_{m-1}x+\beta_{m}},
\]
其中 $n,m$ 为非负整数, $\alpha_{0},\alpha_{1},\cdots,\alpha_{n}$ 与 $\beta_{0},\beta_{1},\cdots,\beta_{m}$ 都为实常数, 且 $\alpha_{0}\neq 0$, $\beta_{0}\neq 0$.

如果 $m>n$, 则称它为真分式; 如果 $m\leqslant n$, 则称它为假分式. 由多项式除法可知,
\[
R(x)=\frac{P(x)}{Q(x)}=\ \text{多项式}\ +\ \text{真分式}.
\]
根据不定积分的简单性质及 $\frac{x^{k+1}}{k+1}$ 为 $x^{k}$ 的原函数可知, 多项式函数的不定积分是容易得到的. 因此, 要计算有理函数的不定积分只需研究真分式的不定积分.

在代数学中, 有部分分式分解定理:

设
\[
R(x)=\frac{P(x)}{Q(x)}
=\frac{\alpha_{0}x^{n}+\alpha_{1}x^{n-1}+\cdots+\alpha_{n-1}x+\alpha_{n}}
{\beta_{0}x^{m}+\beta_{1}x^{m-1}+\cdots+\beta_{m-1}x+\beta_{m}}
\]
为真分式, $\alpha_{0}\neq 0$, $\beta_{0}\neq 0$, $n<m$, 其中
\[
Q(x)=\beta_{0}(x-a_{1})^{\lambda_{1}}\cdots(x-a_{s})^{\lambda_{s}}
(x^{2}+p_{1}x+q_{1})^{\mu_{1}}\cdots(x^{2}+p_{t}x+q_{t})^{\mu_{t}},
\]
这里 $\lambda_{i},\mu_{j}$ ($i=1,2,\cdots,s$; $j=1,2,\cdots,t$) 均为自然数, 且
\[
\sum_{i=1}^{s}\lambda_{i}+2\sum_{j=1}^{t}\mu_{j}=m;\quad p_{j}^{2}-4q_{j}<0,\ j=1,2,\cdots,t.
\]
则 $R(x)$ 可惟一分解为若干部分分式 (称为最简分式) 之和 (即部分分式分解):
\begin{align*}
R(x)&=\frac{P(x)}{Q(x)}
=\frac{A_{1}}{x-a_{1}}+\cdots+\frac{A_{\lambda_{1}}}{(x-a_{1})^{\lambda_{1}}}+\cdots+\frac{B_{1}}{x-a_{s}}\\
&\quad+\cdots+\frac{B_{\lambda_{s}}}{(x-a_{s})^{\lambda_{s}}}
+\frac{K_{1}x+L_{1}}{x^{2}+p_{1}x+q_{1}}+\cdots+\frac{K_{\mu_{1}}x+L_{\mu_{1}}}{(x^{2}+p_{1}x+q_{1})^{\mu_{1}}}\\
&\quad+\cdots+\frac{M_{1}x+N_{1}}{x^{2}+p_{t}x+q_{t}}+\cdots+\frac{M_{\mu_{t}}x+N_{\mu_{t}}}{(x^{2}+p_{t}x+q_{t})^{\mu_{t}}},
\end{align*}
其中 $A_{i},B_{i},K_{i},L_{i},M_{i},N_{i}$ 均为实数.

应用这个定理, 对任何真分式 $R(x)$ 的不定积分, 一旦完成了 $R(x)$ 的部分分式的分解, 最后转化为求各个部分分式 (最简分式) 的不定积分. 由上述讨论知, 任何有理真分式的不定积分都将归为求以下两种形式的不定积分:

(1)
\[
\int\frac{\mathrm{d}x}{(x-a)^{k}}=
\begin{cases}
\ln|x-a|+C, & k=1,\\[2mm]
\dfrac{1}{(1-k)(x-a)^{k-1}}+C, & k>1.
\end{cases}
\]

(2)
\begin{align*}
&\int\frac{Kx+L}{(x^{2}+px+q)^{k}}\mathrm{d}x\quad (p^{2}-4q<0)\\
&\quad=\int\frac{K\left(x+\frac{p}{2}\right)+\left(L-\frac{Kp}{2}\right)}
{\left[\left(x+\frac{p}{2}\right)^{2}+\left(q-\frac{p^{2}}{4}\right)\right]^{k}}\mathrm{d}\left(x+\frac{p}{2}\right)\\
&\quad\xlongequal{t=x+\frac{p}{2}}\int\frac{Kt+\left(L-\frac{Kp}{2}\right)}{(t^{2}+r^{2})^{k}}\mathrm{d}t\\
&\quad=K\int\frac{t}{(t^{2}+r^{2})^{k}}\mathrm{d}t+\left(L-\frac{Kp}{2}\right)\int\frac{\mathrm{d}t}{(t^{2}+r^{2})^{k}}.
\end{align*}

当 $k=1$ 时, 上式右边两个不定积分分别为
\[
\int\frac{t}{t^{2}+r^{2}}\mathrm{d}t=\frac{1}{2}\int\frac{\mathrm{d}(t^{2}+r^{2})}{t^{2}+r^{2}}=\frac{1}{2}\ln(t^{2}+r^{2})+C,
\]
\[
\int\frac{\mathrm{d}t}{t^{2}+r^{2}}=\frac{1}{r}\arctan\frac{t}{r}+C.
\]
当 $k\geqslant 2$ 时, 右边第一个积分为
\begin{align*}
\int\frac{t}{(t^{2}+r^{2})^{k}}\mathrm{d}t&=\frac{1}{2}\int\frac{\mathrm{d}(t^{2}+r^{2})}{(t^{2}+r^{2})^{k}}\\
&=\frac{1}{2(1-k)(t^{2}+r^{2})^{k-1}}+C.
\end{align*}
对于第二个不定积分, 记
\[
I_{k}=\int\frac{\mathrm{d}t}{(t^{2}+r^{2})^{k}},
\]
则
\begin{align*}
I_{k}&=\int\frac{\mathrm{d}t}{(t^{2}+r^{2})^{k}}=\frac{1}{r^{2}}\int\frac{(t^{2}+r^{2})-t^{2}}{(t^{2}+r^{2})^{k}}\mathrm{d}t\\
&=\frac{1}{r^{2}}I_{k-1}-\frac{1}{r^{2}}\int\frac{t^{2}}{(t^{2}+r^{2})^{k}}\mathrm{d}t\\
&=\frac{1}{r^{2}}I_{k-1}+\frac{1}{2r^{2}(k-1)}\int t\,\mathrm{d}\frac{1}{(t^{2}+r^{2})^{k-1}}\\
&\xlongequal{\text{分部积分}}\frac{1}{r^{2}}I_{k-1}+\frac{1}{2r^{2}(k-1)}
\left[\frac{t}{(t^{2}+r^{2})^{k-1}}-I_{k-1}\right],
\end{align*}
\[
I_{k}=\frac{t}{2r^{2}(k-1)(t^{2}+r^{2})^{k-1}}+\frac{2k-3}{2r^{2}(k-1)}I_{k-1}\quad(\text{递推公式}).
\]
或者
\begin{align*}
I_{k-1}&=\int\frac{\mathrm{d}t}{(t^{2}+r^{2})^{k-1}}\\
&\xlongequal{\text{分部积分}}\frac{t}{(t^{2}+r^{2})^{k-1}}-\int t\left[-(k-1)\right]\frac{2t}{(t^{2}+r^{2})^{k}}\mathrm{d}t\\
&=\frac{t}{(t^{2}+r^{2})^{k-1}}+2(k-1)\int\frac{(t^{2}+r^{2})-r^{2}}{(t^{2}+r^{2})^{k}}\mathrm{d}t\\
&=\frac{t}{(t^{2}+r^{2})^{k-1}}+2(k-1)I_{k-1}-2r^{2}(k-1)I_{k},
\end{align*}
\[
I_{k}=\frac{t}{2r^{2}(k-1)(t^{2}+r^{2})^{k-1}}+\frac{2k-3}{2r^{2}(k-1)}I_{k-1}.
\]
根据递推公式, 从 $I_{1}=\int\frac{\mathrm{d}t}{t^{2}+r^{2}}=\frac{1}{r}\arctan\frac{t}{r}+C$ 得到 $I_{2}$, 再从 $I_{2}$ 得到 $I_{3}$, $\cdots$, 从 $I_{k-1}$ 得到 $I_{k}$.

注意, 最后需将 $t=x+\frac{p}{2}$ 代回去.

\begin{example}\label{ex:5.3.1}
求 $\int\frac{x^{2}+1}{(x^{2}-2x+2)^{2}}\mathrm{d}x$.
\end{example}

\begin{solution}
\begin{align*}
\int\frac{x^{2}+1}{(x^{2}-2x+2)^{2}}\mathrm{d}x
&=\int\frac{(x^{2}-2x+2)+(2x-1)}{(x^{2}-2x+2)^{2}}\mathrm{d}x\\
&=\int\frac{\mathrm{d}x}{x^{2}-2x+2}+\int\frac{2x-2}{(x^{2}-2x+2)^{2}}\mathrm{d}x+\int\frac{\mathrm{d}x}{(x^{2}-2x+2)^{2}}\\
&=\int\frac{\mathrm{d}(x-1)}{(x-1)^{2}+1}+\int\frac{\mathrm{d}(x^{2}-2x+2)}{(x^{2}-2x+2)^{2}}+\int\frac{\mathrm{d}(x-1)}{\left[(x-1)^{2}+1\right]^{2}}\\
&=\arctan(x-1)-\frac{1}{x^{2}-2x+2}+\frac{x-1}{2(x^{2}-2x+2)}\\
&\quad+\frac{1}{2}\arctan(x-1)+C\\
&=\frac{3}{2}\arctan(x-1)+\frac{x-3}{2(x^{2}-2x+2)}+C,
\end{align*}
其中
\begin{align*}
I_{2}&=\int\frac{\mathrm{d}t}{(t^{2}+1)^{2}}
\xlongequal[k=2,\ r=1]{\text{递推公式}}\frac{t}{2(t^{2}+1)}+\frac{1}{2}\int\frac{\mathrm{d}t}{t^{2}+1}\\
&=\frac{x-1}{2\left[(x-1)^{2}+1\right]}+\frac{1}{2}\arctan(x-1)+C_{1}\\
&=\frac{x-1}{2(x^{2}-2x+2)}+\frac{1}{2}\arctan(x-1)+C_{1}.
\end{align*}
\end{solution}

\begin{example}\label{ex:5.3.2}
求 $\int\frac{5x+6}{x^{2}+x+1}\mathrm{d}x$.
\end{example}

\begin{solution}
\begin{align*}
\int\frac{5x+6}{x^{2}+x+1}\mathrm{d}x&=\frac{5}{2}\int\frac{2x+\frac{12}{5}}{x^{2}+x+1}\mathrm{d}x
=\frac{5}{2}\int\frac{(2x+1)+\frac{7}{5}}{x^{2}+x+1}\mathrm{d}x\\
&=\frac{5}{2}\int\frac{\mathrm{d}(x^{2}+x+1)}{x^{2}+x+1}
+\frac{7}{2}\int\frac{\mathrm{d}\left(x+\frac{1}{2}\right)}{\left(x+\frac{1}{2}\right)^{2}+\left(\frac{\sqrt{3}}{2}\right)^{2}}\\
&=\frac{5}{2}\ln(x^{2}+x+1)+\frac{7}{2}\cdot\frac{1}{\frac{\sqrt{3}}{2}}\arctan\frac{x+\frac{1}{2}}{\frac{\sqrt{3}}{2}}+C\\
&=\frac{5}{2}\ln(x^{2}+x+1)+\frac{7}{\sqrt{3}}\arctan\frac{2x+1}{\sqrt{3}}+C.
\end{align*}
\end{solution}

\begin{example}\label{ex:5.3.3}
求 $\int\frac{x+1}{x^{2}-4x+3}\mathrm{d}x$.
\end{example}

\begin{solution}[解法 1]
\begin{align*}
\int\frac{x+1}{x^{2}-4x+3}\mathrm{d}x&=\int\frac{x+1}{(x-3)(x-1)}\mathrm{d}x\\
&\xlongequal{\text{凑}}\int\left(\frac{2}{x-3}+\frac{-1}{x-1}\right)\mathrm{d}x\\
&=2\ln|x-3|-\ln|x-1|+C=\ln\left|\frac{(x-3)^{2}}{x-1}\right|+C.
\end{align*}
\end{solution}

\begin{solution}[解法 2 (待定系数法)]
设
\[
\frac{x+1}{x^{2}-4x+3}=\frac{A}{x-3}+\frac{B}{x-1},
\]
则
\[
x+1=A(x-1)+B(x-3)=(A+B)x-(A+3B),\tag{$*$}
\]
比较上面等式两边关于 $x$ 的同幂次的系数知,
\[
\begin{cases}
A+B=1,\\
A+3B=-1.
\end{cases}
\]
下式减上式得 $2B=-2$, $B=-1$, 所以 $A=1-B=2$.

或者, 令 $x=1$ 代入 $(*)$ 式得 $2=-2B$, $B=-1$. 令 $x=3$ 代入 $(*)$ 式得 $4=2A$, $A=2$. 代入即可得到与解法 1 相同的结果.
\end{solution}

\begin{example}\label{ex:5.3.4}
求 $\int\frac{5x+3}{(x^{2}-2x+5)^{2}}\mathrm{d}x$.
\end{example}

\begin{solution}[解法 1]
\begin{align*}
\int&\frac{5x+3}{(x^{2}-2x+5)^{2}}\mathrm{d}x\\
&=\int\frac{\frac{5}{2}(2x-2)}{(x^{2}-2x+5)^{2}}\mathrm{d}x+8\int\frac{\mathrm{d}x}{(x^{2}-2x+5)^{2}}\\
&=\frac{5}{2}\int\frac{\mathrm{d}(x^{2}-2x+5)}{(x^{2}-2x+5)^{2}}+8\int\frac{\mathrm{d}(x-1)}{\left[(x-1)^{2}+2^{2}\right]^{2}}\\
&=-\frac{5}{2}\cdot\frac{1}{x^{2}-2x+5}
+8\left[\frac{1}{8}\cdot\frac{x-1}{(x-1)^{2}+4}+\frac{1}{16}\arctan\frac{x-1}{2}\right]+C\\
&=-\frac{2x-7}{2(x^{2}-2x+5)}+\frac{1}{2}\arctan\frac{x-1}{2}+C.
\end{align*}
\end{solution}

\begin{solution}[解法 2]
\begin{align*}
\int\frac{5x+3}{(x^{2}-2x+5)^{2}}\mathrm{d}x
&\xlongequal{u=x-1}\int\frac{5u+8}{(u^{2}+4)^{2}}\mathrm{d}u\\
&=\frac{5}{2}\int\frac{\mathrm{d}(u^{2}+4)}{(u^{2}+4)^{2}}+2\int\frac{(u^{2}+4)-u^{2}}{(u^{2}+4)^{2}}\mathrm{d}u\\
&=-\frac{5}{2}\cdot\frac{1}{u^{2}+4}+2\int\frac{\mathrm{d}u}{u^{2}+4}+\int u\,\mathrm{d}\frac{1}{u^{2}+4}\\
&=-\frac{5}{2(u^{2}+4)}+2\int\frac{\mathrm{d}u}{u^{2}+4}+\frac{u}{u^{2}+4}-\int\frac{\mathrm{d}u}{u^{2}+4}\\
&=\frac{2u-5}{2(u^{2}+4)}+\frac{1}{2}\arctan\frac{u}{2}+C\\
&=\frac{2x-7}{2(x^{2}-2x+5)}+\frac{1}{2}\arctan\frac{x-1}{2}+C.
\end{align*}
\end{solution}

\begin{example}\label{ex:5.3.5}
求 $\int\frac{x^{5}+x^{3}}{x^{3}+1}\mathrm{d}x$.
\end{example}

\begin{solution}
\begin{align*}
\int\frac{x^{5}+x^{3}}{x^{3}+1}\mathrm{d}x
&=\int\left[(x^{2}+1)-\frac{x^{2}+1}{x^{3}+1}\right]\mathrm{d}x\\
&=\int\left[(x^{2}+1)-\frac{2}{3}\cdot\frac{1}{x+1}-\frac{1}{3}\cdot\frac{x+1}{x^{2}-x+1}\right]\mathrm{d}x\\
&=\int\left(x^{2}+1-\frac{2}{3}\cdot\frac{1}{x+1}-\frac{1}{6}\cdot\frac{2x-1}{x^{2}-x+1}-\frac{1}{2}\cdot\frac{1}{x^{2}-x+1}\right)\mathrm{d}x\\
&=\frac{x^{3}}{3}+x-\frac{1}{6}\ln(x^{3}+1)(x+1)^{3}-\frac{1}{\sqrt{3}}\arctan\frac{2(x-1)}{\sqrt{3}}+C.
\end{align*}
其中 $\frac{x^{5}+x^{3}}{x^{3}+1}=(x^{2}+1)-\frac{x^{2}+1}{x^{3}+1}$ 可用凑的方法, 也可用多项式的辗转相除得到. 而应用待定系数法可令
\[
\frac{x^{2}+1}{x^{3}+1}=\frac{A}{x+1}+\frac{Bx+C}{x^{2}-x+1},
\]
则有
\begin{align*}
x^{2}+1&=A(x^{2}-x+1)+(Bx+C)(x+1)\\
&=(A+B)x^{2}+(-A+B+C)x+(A+C),
\end{align*}
从而有
\[
\left\{\begin{array}
{l@{\qquad}r}
A+B=1, & \text{(1)}\\
-A+B+C=0, & \text{(2)}\\
A+C=1. & \text{(3)}
\end{array}\right.
\]
于是, 由 (1)$+$(2) 及 (2)$+$(3) 得到
\[
\left\{\begin{array}
{l@{\qquad}r}
2B+C=1, & \text{(4)}\\
B+2C=1. & \text{(5)}
\end{array}\right.
\]
再由 $2\times$(5)$-$(4) 得 $3C=1$, $C=\frac{1}{3}$; $B=1-2C=1-\frac{2}{3}=\frac{1}{3}$; $A=1-B=1-\frac{1}{3}=\frac{2}{3}$.
\end{solution}

\begin{example}\label{ex:5.3.6}
求 $\int\frac{x^{3}+1}{x^{4}-3x^{3}+3x^{2}-x}\mathrm{d}x$.
\end{example}

\begin{solution}[解法 1]
$x^{4}-3x^{3}+3x^{2}-x=x(x-1)^{3}$. 用待定系数法, 令
\begin{align*}
\frac{x^{3}+1}{x(x-1)^{3}}&=\frac{x^{3}+1}{x^{4}-3x^{3}+3x^{2}-x}\\
&=\frac{A}{x}+\frac{B}{(x-1)^{3}}+\frac{C}{(x-1)^{2}}+\frac{D}{x-1},
\end{align*}
则
\begin{align*}
x^{3}+1&=A(x-1)^{3}+Bx+Cx(x-1)+Dx(x-1)^{2}\\
&=(A+D)x^{3}+(-3A+C-2D)x^{2}+(3A+B-C+D)x-A,
\end{align*}
于是
\[
\begin{cases}
A+D=1,\\
-3A+C-2D=0,\\
3A+B-C+D=0,\\
-A=1.
\end{cases}\quad
\text{即}\quad
\begin{cases}
A=-1,\\
D=1-A=1-(-1)=2,\\
C=3A+2D=-3+4=1,\\
B=-3A+C-D=3+1-2=2.
\end{cases}
\]
所以
\begin{align*}
&\int\frac{x^{3}+1}{x^{4}-3x^{3}+3x^{2}-x}\mathrm{d}x\\
&=\int\left[\frac{-1}{x}+\frac{2}{(x-1)^{3}}+\frac{1}{(x-1)^{2}}+\frac{2}{x-1}\right]\mathrm{d}x\\
&=-1\cdot\ln|x|-\frac{1}{(x-1)^{2}}-\frac{1}{x-1}+2\ln|x-1|+C\\
&=\ln\frac{(x-1)^{2}}{|x|}-\frac{x}{(x-1)^{2}}+C.
\end{align*}
\end{solution}

\begin{solution}[解法 2]
设 $x^{3}+1=A(x-1)^{3}+Bx+Cx(x-1)+Dx(x-1)^{2}$. 令 $x=1$ 代入上式得 $B=2$; $x=0$ 代入上式得 $A=-1$. 从而
\[
x^{3}+1=-(x-1)^{3}+2x+Cx(x-1)+Dx(x-1)^{2},
\]
\begin{align*}
Cx(x-1)+Dx(x-1)^{2}&=x^{3}+1+x^{3}-3x^{2}+3x-1-2x\\
&=2x^{3}-3x^{2}+x=x(2x^{2}-3x+1)\\
&=x(x-1)(2x-1).
\end{align*}
于是, 当 $x\neq 0$, $x\neq 1$ 时,
\[
C+D(x-1)=2x-1.
\]
令 $x\to 1$ 得到 $C=1$; 令 $x\to 0$ 得到 $1-D=C-D=-1$, $D=2$.

综上所述有 $A=-1$, $B=2$, $C=1$, $D=2$. 积分过程同解法 1.
\end{solution}

\begin{example}\label{ex:5.3.7}
求 $\int\frac{\mathrm{d}x}{x^{4}+1}$.
\end{example}

\begin{solution}[解法 1]
\begin{align*}
\int\frac{\mathrm{d}x}{x^{4}+1}&=\int\frac{\mathrm{d}x}{(x^{2}+1)^{2}-(\sqrt{2}x)^{2}}
=\int\frac{\mathrm{d}x}{(x^{2}+\sqrt{2}x+1)(x^{2}-\sqrt{2}x+1)}\\
&=\int\left(\frac{\frac{\sqrt{2}}{4}x+\frac{1}{2}}{x^{2}+\sqrt{2}x+1}+\frac{-\frac{\sqrt{2}}{4}x+\frac{1}{2}}{x^{2}-\sqrt{2}x+1}\right)\mathrm{d}x\\
&=\frac{1}{4\sqrt{2}}\int\left(\frac{2x+\sqrt{2}}{x^{2}+\sqrt{2}x+1}-\frac{2x-\sqrt{2}}{x^{2}-\sqrt{2}x+1}\right)\mathrm{d}x\\
&\quad+\frac{1}{4}\int\left[\frac{1}{\left(x+\frac{\sqrt{2}}{2}\right)^{2}+\left(\frac{1}{\sqrt{2}}\right)^{2}}
+\frac{1}{\left(x-\frac{\sqrt{2}}{2}\right)^{2}+\left(\frac{1}{\sqrt{2}}\right)^{2}}\right]\mathrm{d}x\\
&=\frac{1}{4\sqrt{2}}\ln\frac{x^{2}+\sqrt{2}x+1}{x^{2}-\sqrt{2}x+1}
+\frac{1}{2\sqrt{2}}\left[\arctan(\sqrt{2}x+1)+\arctan(\sqrt{2}x-1)\right]+C.
\end{align*}
\end{solution}

\begin{solution}[解法 2]
因为
\begin{align*}
\int\frac{x^{2}+1}{x^{4}+1}\mathrm{d}x&=\int\frac{1+\frac{1}{x^{2}}}{x^{2}+\frac{1}{x^{2}}}\mathrm{d}x
=\int\frac{\mathrm{d}\left(x-\frac{1}{x}\right)}{\left(x-\frac{1}{x}\right)^{2}+2}\\
&=\frac{1}{\sqrt{2}}\arctan\frac{x-\frac{1}{x}}{\sqrt{2}}+C_{1},
\end{align*}
\begin{align*}
\int\frac{1-x^{2}}{x^{4}+1}\mathrm{d}x&=\int\frac{\frac{1}{x^{2}}-1}{x^{2}+\frac{1}{x^{2}}}\mathrm{d}x
=-\int\frac{\mathrm{d}\left(x+\frac{1}{x}\right)}{\left(x+\frac{1}{x}\right)^{2}-2}\\
&=-\frac{1}{2\sqrt{2}}\int\left[\frac{1}{\left(x+\frac{1}{x}\right)-\sqrt{2}}
-\frac{1}{\left(x+\frac{1}{x}\right)+\sqrt{2}}\right]\mathrm{d}\left(x+\frac{1}{x}\right)\\
&=\frac{1}{2\sqrt{2}}\ln\frac{x+\frac{1}{x}+\sqrt{2}}{x+\frac{1}{x}-\sqrt{2}}+C_{2}
=\frac{1}{2\sqrt{2}}\ln\frac{x^{2}+\sqrt{2}x+1}{x^{2}-\sqrt{2}x+1}+C_{2},
\end{align*}
所以
\begin{align*}
\int\frac{\mathrm{d}x}{x^{4}+1}&=\frac{1}{2}\int\left(\frac{1-x^{2}}{x^{4}+1}+\frac{x^{2}+1}{x^{4}+1}\right)\mathrm{d}x\\
&=\frac{1}{4\sqrt{2}}\ln\frac{x^{2}+\sqrt{2}x+1}{x^{2}-\sqrt{2}x+1}+\frac{1}{2\sqrt{2}}\arctan\frac{x-\frac{1}{x}}{\sqrt{2}}+C_{3}\\
&=\frac{1}{4\sqrt{2}}\ln\frac{x^{2}+\sqrt{2}x+1}{x^{2}-\sqrt{2}x+1}
+\frac{1}{2\sqrt{2}}\left[\arctan(\sqrt{2}x+1)+\arctan(\sqrt{2}x-1)\right]+C,
\end{align*}
其中, 因为
\begin{align*}
\tan\left[\arctan(\sqrt{2}x+1)+\arctan(\sqrt{2}x-1)\right]
&=\frac{(\sqrt{2}x+1)+(\sqrt{2}x-1)}{1-(2x^{2}-1)}=\frac{\sqrt{2}x}{1-x^{2}},\\
\tan\left(\arctan\frac{x-\frac{1}{x}}{\sqrt{2}}\right)&=\frac{x-\frac{1}{x}}{\sqrt{2}}=\frac{x^{2}-1}{\sqrt{2}x}\\
&=-\cot\left[\arctan(\sqrt{2}x+1)+\arctan(\sqrt{2}x-1)\right],
\end{align*}
所以
\[
\arctan\frac{x-\frac{1}{x}}{\sqrt{2}}=\frac{\pi}{2}+\left[\arctan(\sqrt{2}x+1)+\arctan(\sqrt{2}x-1)\right].
\]
\end{solution}

代数学关于有理函数的部分分式分解定理从理论上证明了任何有理函数的不定积分都为初等函数. 但该有理函数的分母 $Q(x)$, 当它的次数很高时, 具体分解因子是十分困难的. 即使 $Q(x)$ 能因式分解, 用待定系数法将此有理函数化为最简分式之和的计算也会十分烦琐.

现在来介绍几类函数的不定积分, 它们可化为有理函数的不定积分. 因此, 理论上这些函数的不定积分是可以积出来的.

由 $u(x)$, $v(x)$ 及常数经过有限次四则运算所得到的函数称为关于 $u(x)$, $v(x)$ 的有理式, 并表示为 $R(u(x),v(x))=\frac{P(u(x),v(x))}{Q(u(x),v(x))}$, 其中 $P(x,y)=\sum\limits_{i=1}^{n}\sum\limits_{j=1}^{m}a_{ij}x^{i}y^{j}$ ($a_{ij}$ 为实系数, $x,y\in\mathbb{R}$) 为实二元多项式, $Q(x,y)$ 也为实二元多项式.

\subsection*{1. 三角函数有理式的不定积分}

\[
\int R(\sin x,\cos x)\mathrm{d}x
\xlongequal{t=\tan\frac{x}{2}}\int R\left(\frac{2t}{1+t^{2}},\frac{1-t^{2}}{1+t^{2}}\right)\frac{2}{1+t^{2}}\mathrm{d}t.
\]
等式右端是 $t$ 的有理函数的不定积分, 因此, 理论上它可以被积出来, 成为一个初等函数.

正因为如此, 我们称 $t=\tan\frac{x}{2}$ 为万能变换. 其中
\begin{gather*}
x=2\arctan t,\quad \mathrm{d}x=\frac{2}{1+t^{2}}\mathrm{d}t,\\
\sin x=\frac{2\sin\frac{x}{2}\cos\frac{x}{2}}{\sin^{2}\frac{x}{2}+\cos^{2}\frac{x}{2}}
=\frac{2\tan\frac{x}{2}}{\tan^{2}\frac{x}{2}+1}=\frac{2t}{1+t^{2}},\\
\cos x=\frac{\cos^{2}\frac{x}{2}-\sin^{2}\frac{x}{2}}{\sin^{2}\frac{x}{2}+\cos^{2}\frac{x}{2}}
=\frac{1-\tan^{2}\frac{x}{2}}{1+\tan^{2}\frac{x}{2}}=\frac{1-t^{2}}{1+t^{2}}.
\end{gather*}

上面所用的万能变换 $t=\tan\frac{x}{2}$ 对三角函数有理式的不定积分理论上总是有效的. 但如果右边积分中被积函数的分母的次数太高, 往往难以积出来. 此时 ``万能'' 就成为一个虚名了.

另外, 还有几个可化为有理函数的不定积分的变换, 有时这些有理函数分母的幂次较低, 积分比较容易.

(1) 若 $R(-\sin x,\cos x)=-R(\sin x,\cos x)$, 令 $t=\cos x$;

(2) 若 $R(\sin x,-\cos x)=-R(\sin x,\cos x)$, 令 $t=\sin x$;

(3) 若 $R(-\sin x,-\cos x)=R(\sin x,\cos x)$, 令 $t=\tan x$.

\begin{example}\label{ex:5.3.8}
求 $\int\frac{1+\sin x}{\sin x(1+\cos x)}\mathrm{d}x$.
\end{example}

\begin{solution}
\begin{align*}
&\int\frac{1+\sin x}{\sin x(1+\cos x)}\mathrm{d}x\\
&\xlongequal{t=\tan\frac{x}{2}}\int\frac{1+\frac{2t}{1+t^{2}}}{\frac{2t}{1+t^{2}}\left(1+\frac{1-t^{2}}{1+t^{2}}\right)}\cdot\frac{2}{1+t^{2}}\mathrm{d}t\\
&=\int\frac{1+t^{2}+2t}{t(1+t^{2}+1-t^{2})}\mathrm{d}t=\frac{1}{2}\int\left(t+2+\frac{1}{t}\right)\mathrm{d}t\\
&=\frac{1}{2}\left(\frac{t^{2}}{2}+2t+\ln|t|\right)+C\\
&=\frac{1}{4}\tan^{2}\frac{x}{2}+\tan\frac{x}{2}+\frac{1}{2}\ln\left|\tan\frac{x}{2}\right|+C.
\end{align*}
\end{solution}

\begin{example}\label{ex:5.3.9}
求 $\int\frac{\mathrm{d}x}{a^{2}\sin^{2}x+b^{2}\cos^{2}x}$ ($a>0$, $b>0$).
\end{example}

\begin{solution}
\begin{align*}
\int\frac{\mathrm{d}x}{a^{2}\sin^{2}x+b^{2}\cos^{2}x}
&=\int\frac{\frac{1}{\cos^{2}x}}{a^{2}\frac{\sin^{2}x}{\cos^{2}x}+b^{2}}\mathrm{d}x\\
&=\int\frac{\mathrm{d}\tan x}{a^{2}\tan^{2}x+b^{2}}
=\frac{1}{ab}\arctan\left(\frac{a}{b}\arctan x\right)+C.
\end{align*}
\end{solution}

\begin{example}\label{ex:5.3.10}
求 $\int\frac{\mathrm{d}x}{\sin^{2}x\cos x}$.
\end{example}

\begin{solution}
\begin{align*}
\int\frac{\mathrm{d}x}{\sin^{2}x\cos x}&=\int\frac{\mathrm{d}\sin x}{\sin^{2}x(1-\sin^{2}x)}
\xlongequal{t=\sin x}\int\frac{\mathrm{d}t}{t^{2}(1-t^{2})}\\
&=\int\left(\frac{1}{t^{2}}+\frac{1}{1-t^{2}}\right)\mathrm{d}t=-\frac{1}{t}+\frac{1}{2}\ln\frac{1+t}{1-t}+C\\
&=-\frac{1}{\sin x}+\frac{1}{2}\ln\frac{1+\sin x}{1-\sin x}+C.
\end{align*}
\end{solution}

\begin{example}\label{ex:5.3.11}
求 $\int\frac{\sin^{4}x}{\cos^{2}x}\mathrm{d}x$.
\end{example}

\begin{solution}[解法 1]
\begin{align*}
\int\frac{\sin^{4}x}{\cos^{2}x}\mathrm{d}x&=\int\frac{\sin^{4}x}{(\sin^{2}x+\cos^{2}x)^{2}}\mathrm{d}\tan x
=\int\frac{\tan^{4}x}{(\tan^{2}x+1)^{2}}\mathrm{d}\tan x\\
&\xlongequal{t=\tan x}\int\frac{t^{4}}{(t^{2}+1)^{2}}\mathrm{d}t
=\int\frac{(t^{4}+2t^{2}+1)-2(t^{2}+1)+1}{(t^{2}+1)^{2}}\mathrm{d}t\\
&=\int\left[1-\frac{1}{1+t^{2}}+\frac{1}{(1+t^{2})^{2}}\right]\mathrm{d}t\\
&\xlongequal{\text{递推公式}}t-2\arctan t+\left(\frac{1}{2}\cdot\frac{t}{1+t^{2}}+\frac{1}{2}\arctan t\right)+C\\
&=t-\frac{3}{2}\arctan t+\frac{t}{2(1+t^{2})}+C\\
&=\tan x-\frac{3}{2}x+\frac{1}{2}\cdot\frac{\tan x}{\sec^{2}x}+C\\
&=\tan x-\frac{3}{2}x+\frac{1}{4}\sin 2x+C.
\end{align*}
\end{solution}

\begin{solution}[解法 2]
\begin{align*}
\int\frac{\sin^{4}x}{\cos^{2}x}\mathrm{d}x&\xlongequal{t=\tan x}\int\frac{t^{4}}{(t^{2}+1)^{2}}\mathrm{d}t\\
&=\int\left[1-\frac{1}{t^{2}+1}-\frac{t^{2}}{(t^{2}+1)^{2}}\right]\mathrm{d}t\\
&=t-\arctan t+\frac{1}{2}\int t\,\mathrm{d}\frac{1}{t^{2}+1}\\
&=t-\arctan t+\frac{t}{2(t^{2}+1)}-\frac{1}{2}\int\frac{\mathrm{d}t}{t^{2}+1}\\
&=t-\frac{3}{2}\arctan t+\frac{t}{2(t^{2}+1)}+C\\
&=\tan x-\frac{3}{2}x+\frac{1}{2}\cdot\frac{\tan x}{\sec^{2}x}+C\\
&=\tan x-\frac{3}{2}x+\frac{1}{4}\sin 2x+C.
\end{align*}
\end{solution}

\begin{example}\label{ex:5.3.12}
求 Poisson 积分 $\int\frac{1-r^{2}}{1-2r\cos x+r^{2}}\mathrm{d}x$, $0<r<1$, $|x|<\pi$.
\end{example}

\begin{solution}
作万能变换 $t=\tan\frac{x}{2}$, 得到
\begin{align*}
&\int\frac{1-r^{2}}{1-2r\cos x+r^{2}}\mathrm{d}x\\
&=\int\frac{1-r^{2}}{1-2r\frac{1-t^{2}}{1+t^{2}}+r^{2}}\cdot\frac{2}{1+t^{2}}\mathrm{d}t\\
&=2(1-r^{2})\int\frac{\mathrm{d}t}{1+t^{2}-2r+2rt^{2}+r^{2}+r^{2}t^{2}}\\
&=2(1-r^{2})\int\frac{\mathrm{d}t}{(1-r)^{2}+(1+r)^{2}t^{2}}\\
&=\frac{2(1-r^{2})}{(1+r)^{2}}\int\frac{\mathrm{d}t}{\left(\frac{1-r}{1+r}\right)^{2}+t^{2}}\\
&=\frac{2(1-r^{2})}{(1+r)^{2}}\cdot\frac{1}{\frac{1-r}{1+r}}\arctan\frac{t}{\frac{1-r}{1+r}}+C\\
&=2\arctan\left(\frac{1+r}{1-r}\tan\frac{x}{2}\right)+C.
\end{align*}
\end{solution}

\subsection*{2. 某些带无理根式的不定积分}

设 $R(x,y_{1},\cdots,y_{k})$ 为 $k+1$ 元实有理函数.

(1) $\int R\left(x,\sqrt[n]{\frac{ax+b}{cx+d}},\cdots,\sqrt[m]{\frac{ax+b}{cx+d}}\right)\mathrm{d}x$ 型不定积分, 其中 $a,b,c,d\in\mathbb{R}$, $ad-bc\neq 0$, $n,\cdots,m\in\mathbb{N}$. 令
\[
t=\left(\frac{ax+b}{cx+d}\right)^{\frac{1}{s}},
\]
这里 $s$ 为 $n,\cdots,m$ 的最小公倍数, 则
\begin{gather*}
\frac{ax+b}{cx+d}=t^{s},\\
ax+b=ct^{s}x+t^{s}d,\\
x=\frac{dt^{s}-b}{a-ct^{s}},\\
\mathrm{d}x=\frac{st^{s-1}(ad-bc)}{(a-ct^{s})^{2}}\mathrm{d}t.
\end{gather*}

(2) $\int R(x,\sqrt{ax^{2}+bx+c})\mathrm{d}x$ 型不定积分, 其中当 $a>0$ 时, $b^{2}-4ac\neq 0$; 当 $a<0$ 时, $b^{2}-4ac>0$. 由于
\[
ax^{2}+bx+c=a\left[\left(x+\frac{b}{2a}\right)^{2}+\frac{4ac-b^{2}}{4a^{2}}\right],
\]
若记 $u=x+\frac{b}{2a}$, $k^{2}=\left|\frac{4ac-b^{2}}{4a^{2}}\right|$, 则上述二次三项式必属于以下三种情形之一:
\[
|a|(u^{2}+k^{2}),\quad |a|(u^{2}-k^{2}),\quad |a|(k^{2}-u^{2}).
\]
因此, 上面的不定积分转化为下面三种类型不定积分之一:
\[
\int R(u,\sqrt{u^{2}\pm k^{2}})\mathrm{d}u,\quad \int R(u,\sqrt{k^{2}-u^{2}})\mathrm{d}u.
\]
当分别令 $u=k\tan t$, $u=k\sec t$, $u=k\sin t$ 后, 它们都可化为三角有理式的不定积分. 特别地, 还有下面几种情况.

① $a>0$, 作变换 $\sqrt{ax^{2}+bx+c}=t-\sqrt{a}x$ (或 $\sqrt{a}x\pm t$), 于是,
\begin{gather*}
t=\sqrt{ax^{2}+bx+c}+\sqrt{a}x,\\
ax^{2}+bx+c=t^{2}-2\sqrt{a}tx+ax^{2},
\end{gather*}
两边消去 $ax^{2}$, 移项得
\begin{gather*}
(b+2\sqrt{a}t)x=t^{2}-c,\\
x=\frac{t^{2}-c}{b+2\sqrt{a}t},\\
\mathrm{d}x=\frac{2t(b+2\sqrt{a}t)-2\sqrt{a}(t^{2}-c)}{(b+2\sqrt{a}t)^{2}}\mathrm{d}t\\
=\frac{2bt+2\sqrt{a}t^{2}+2\sqrt{ac}}{(b+2\sqrt{a}t)^{2}}\mathrm{d}t.
\end{gather*}

② $c>0$, 作变换 $\sqrt{ax^{2}+bx+c}=tx+\sqrt{c}$ (或 $tx-\sqrt{c}$), 则
\begin{gather*}
t=\frac{\sqrt{ax^{2}+bx+c}-\sqrt{c}}{x},\\
ax^{2}+bx+c=t^{2}x^{2}+2\sqrt{c}tx+c.
\end{gather*}
两边消去 $c$, 约去 $x$ 得
\begin{gather*}
ax+b=t^{2}x+2\sqrt{c}t,\\
x=\frac{2\sqrt{c}t-b}{a-t^{2}},\\
\mathrm{d}x=\frac{2a\sqrt{c}-2bt+2\sqrt{c}t^{2}}{(a-t^{2})^{2}}\mathrm{d}t.
\end{gather*}

③ $b^{2}-4ac>0$, 则 $ax^{2}+bx+c=0$ 有两个相异的实根 $\lambda,\mu$. 于是
\[
ax^{2}+bx+c=a(x-\lambda)(x-\mu).
\]
作变换
\[
\sqrt{ax^{2}+bx+c}=t(x-\lambda),
\]
两边平方得
\begin{gather*}
a(x-\lambda)(x-\mu)=ax^{2}+bx+c=t^{2}(x-\lambda)^{2},\\
a(x-\mu)=t^{2}(x-\lambda),\\
x=\frac{a\mu-\lambda t^{2}}{a-t^{2}},\\
\mathrm{d}x=\frac{2a(\mu-\lambda)t}{(a-t^{2})^{2}}\mathrm{d}t.
\end{gather*}

④ $b^{2}-4ac=0$, 当 $a>0$, 则 $\sqrt{ax^{2}+bx+c}=\sqrt{a(x-\lambda)^{2}}=\sqrt{a}(x-\lambda)$ 或 $\sqrt{a}(\lambda-x)$, 此时, $R(x,\sqrt{ax^{2}+bx+c})$ 为 $x$ 的有理式;

当 $a<0$ 时, $\sqrt{ax^{2}+bx+c}=\sqrt{a(x-\lambda)^{2}}$ 对 $x\neq\lambda$ 无意义.

⑤ $b^{2}-4ac<0$, 则 $ax^{2}+bx+c=a\left[\left(x+\frac{b}{2a}\right)^{2}+\frac{4ac-b^{2}}{4a^{2}}\right]$ 与 $a$ 同号.

当 $a>0$ 时, 由 ① 作变换 $x=\frac{t^{2}-c}{b+2\sqrt{a}t}$;

当 $a<0$ 时, $\sqrt{ax^{2}+bx+c}$ 无意义.

\begin{example}\label{ex:5.3.13}
求 $\int\frac{\mathrm{d}x}{x\sqrt{x^{2}-2x-3}}$.
\end{example}

\begin{solution}[解法 1]
\begin{align*}
\int\frac{\mathrm{d}x}{x\sqrt{x^{2}-2x-3}}&=\int\frac{\mathrm{d}(x-1)}{x\sqrt{(x-1)^{2}-4}}\\
&\xlongequal{u=x-1}\int\frac{\mathrm{d}u}{(u+1)\sqrt{u^{2}-4}}
\xlongequal{u=2\sec\theta}\int\frac{2\sec\theta\cdot\tan\theta}{(2\sec\theta+1)\cdot 2\tan\theta}\mathrm{d}\theta\\
&=\int\frac{\mathrm{d}\theta}{2+\cos\theta}
\xlongequal{t=\tan\frac{\theta}{2}}\int\frac{\frac{2}{1+t^{2}}}{2+\frac{1-t^{2}}{1+t^{2}}}\mathrm{d}t
=\int\frac{2}{t^{2}+3}\mathrm{d}t=\frac{2}{\sqrt{3}}\arctan\frac{t}{\sqrt{3}}+C\\
&=\frac{2}{\sqrt{3}}\arctan\left(\frac{1}{\sqrt{3}}\tan\frac{\theta}{2}\right)+C
=\frac{2}{\sqrt{3}}\arctan\frac{\sqrt{x^{2}-2x-3}}{\sqrt{3}(x+1)}+C.
\end{align*}
其中
\begin{align*}
\tan\frac{\theta}{2}&=\frac{2\sin\frac{\theta}{2}\cos\frac{\theta}{2}}{2\cos^{2}\frac{\theta}{2}}
=\frac{\sin\theta}{1+\cos\theta}=\frac{\tan\theta}{\sec\theta+1}\\
&=\frac{\sqrt{\left(\frac{u}{2}\right)^{2}-1}}{\frac{u}{2}+1}
=\frac{\sqrt{\left(\frac{x-1}{2}\right)^{2}-1}}{\frac{x-1}{2}+1}
=\frac{\sqrt{x^{2}-2x-3}}{x+1}.
\end{align*}
\end{solution}

\begin{solution}[解法 2]
$a=1>0$, 令 $\sqrt{x^{2}-2x-3}=x-t$, 两边平方得到
\begin{gather*}
x^{2}-2x-3=x^{2}-2tx+t^{2},\\
-2x-3=-2tx+t^{2},\\
x=\frac{t^{2}+3}{2(t-1)},\\
\mathrm{d}x=\frac{2t(t-1)-(t^{2}+3)}{2(t-1)^{2}}\mathrm{d}t=\frac{t^{2}-2t-3}{2(t-1)^{2}}\mathrm{d}t,\\
\sqrt{x^{2}-2x-3}=x-t=\frac{t^{2}+3}{2(t-1)}-t=\frac{-(t^{2}-2t-3)}{2(t-1)}.
\end{gather*}
于是,
\begin{align*}
\int\frac{\mathrm{d}x}{x\sqrt{x^{2}-2x-3}}
&=\int\frac{\frac{t^{2}-2t-3}{2(t-1)^{2}}}{\frac{t^{2}+3}{2(t-1)}\cdot\frac{-(t^{2}-2t-3)}{2(t-1)}}\mathrm{d}t\\
&=-2\int\frac{\mathrm{d}t}{t^{2}+3}=-\frac{2}{\sqrt{3}}\arctan\frac{t}{\sqrt{3}}+C_{1}\\
&=\frac{2}{\sqrt{3}}\arctan\frac{\sqrt{x^{2}-2x-3}-x}{\sqrt{3}}+C_{1}\\
&=\frac{2}{\sqrt{3}}\left(\arctan\frac{\sqrt{x^{2}-2x-3}}{\sqrt{3}(x+1)}-\frac{\pi}{3}\right)+C_{1}\\
&=\frac{2}{\sqrt{3}}\arctan\frac{\sqrt{x^{2}-2x-3}}{\sqrt{3}(x+1)}+C.
\end{align*}
\end{solution}

读者必须知道, 存在初等函数, 它的不定积分 ``积不出来'', 即原函数不能用初等函数表达. 例如:
\[
\int\mathrm{e}^{\pm x^{2}}\mathrm{d}x,\quad \int\sin x^{2}\,\mathrm{d}x,\quad \int\cos x^{2}\,\mathrm{d}x,\quad
\mathrm{six}=\int\frac{\sin x}{x}\mathrm{d}x\ (\text{积分正弦}),
\]
\[
\mathrm{cix}=\int\frac{\cos x}{x}\mathrm{d}x\ (\text{积分余弦}),\quad
\mathrm{Lix}=\int\frac{\mathrm{d}x}{\ln x}\ (\text{积分对数}),
\]
Legendre 形式下的第一、二、三类椭圆积分 ($0<k<1$)
\[
\int\frac{\mathrm{d}x}{\sqrt{1-k^{2}\sin^{2}x}},\quad \int\sqrt{1-k^{2}\sin^{2}x}\,\mathrm{d}x,\quad
\int\frac{\mathrm{d}x}{(1+h\sin^{2}x)\sqrt{1-k^{2}\sin^{2}x}}.
\]
虽然它们都存在 (参阅第 6 章微积分基本定理~\ref{thm:6.3.2}), 但不能用初等函数表达 (证明非常艰难, Liouville 于 1835 年作出证明). 因此, 初等函数的原函数不一定是初等函数.

\begin{exblock}{练习题 5.3}
  \item 求下列有理函数的不定积分:
  \begin{enumerate}[label=(\arabic*)]
    \item $\int\frac{x^{3}}{x-1}\mathrm{d}x$;
    \item $\int\frac{x-2}{x^{2}-7x+12}\mathrm{d}x$;
    \item $\int\frac{\mathrm{d}x}{1+x^{3}}$;
    \item $\int\frac{\mathrm{d}x}{1+x^{4}}$;
    \item $\int\frac{\mathrm{d}x}{(x-1)(x^{2}+1)^{2}}$;
    \item $\int\frac{x-2}{(2x^{2}+2x+1)^{2}}\mathrm{d}x$;
    \item $\int\frac{5x+6}{x^{2}+x+1}\mathrm{d}x$;
    \item $\int\frac{x}{2x^{2}-3x-2}\mathrm{d}x$;
    \item $\int\frac{2x^{2}+1}{(x+3)(x-1)(x-4)}\mathrm{d}x$;
    \item $\int\frac{\mathrm{d}x}{2x^{3}+3x^{2}+x}$;
    \item $\int\frac{x^{2}}{x^{3}+5x^{2}+8x+4}\mathrm{d}x$;
    \item $\int\frac{4x+3}{(x-2)^{3}}\mathrm{d}x$;
    \item $\int\frac{x^{2}}{(x+2)^{2}(x+4)^{2}}\mathrm{d}x$;
    \item $\int\frac{8x^{3}+7}{(x+1)(2x+1)^{3}}\mathrm{d}x$;
    \item $\int\frac{\mathrm{d}x}{(x+1)(x+2)^{2}(x+3)^{3}}$;
    \item $\int\frac{\mathrm{d}x}{(x^{2}+9)^{3}}$;
    \item $\int\frac{1+x^{2}}{1+x^{4}}\mathrm{d}x$;
    \item $\int\frac{\mathrm{d}x}{x^{4}(1+x^{2})}$;
    \item $\int\frac{x-5}{x^{3}-3x^{2}+4}\mathrm{d}x$;
    \item $\int\frac{x^{7}}{x^{4}+2}\mathrm{d}x$.
  \end{enumerate}

  \item 求下列不定积分:
  \begin{enumerate}[label=(\arabic*)]
    \item $\int\frac{\mathrm{d}x}{\sqrt{x^{2}+x}}$;
    \item $\int\frac{1}{x^{2}}\sqrt{\frac{1-x}{1+x}}\mathrm{d}x$;
    \item $\int\frac{\mathrm{d}x}{\sqrt{x}(1+\sqrt[3]{x})}$;
    \item $\int\sqrt{\frac{x-a}{b-x}}\mathrm{d}x$, $a<x<b$;
    \item $\int\frac{\mathrm{d}x}{1+\sqrt{x}}$;
    \item $\int\frac{\mathrm{d}x}{x\sqrt{x^{2}-1}}$;
    \item $\int\frac{\mathrm{d}x}{1+\sqrt[3]{1+x}}$;
    \item $\int\frac{\sqrt[3]{x}}{x(\sqrt{x}+\sqrt[3]{x})}\mathrm{d}x$;
    \item $\int\frac{x}{\sqrt{x+1}+\sqrt[3]{x+1}}\mathrm{d}x$;
    \item $\int\frac{\sqrt{1+x^{2}}}{2+x^{2}}\mathrm{d}x$;
    \item $\int\frac{\mathrm{d}x}{x\sqrt{1+x^{2}}}$;
    \item $\int\frac{x^{3}}{\sqrt{1+2x^{2}}}\mathrm{d}x$;
    \item $\int\sqrt{\frac{1-\sqrt{x}}{1+\sqrt{x}}}\mathrm{d}x$;
    \item $\int\frac{\mathrm{d}x}{x\sqrt{a^{2}-x^{2}}}$, $|x|<a$, $x\neq 0$;
    \item $\int\frac{x^{2}}{\sqrt{1+x-x^{2}}}\mathrm{d}x$;
    \item $\int\frac{\mathrm{d}x}{(x+a)^{2}(x+b)^{3}}$.
  \end{enumerate}

  \item 求下列不定积分:
  \begin{enumerate}[label=(\arabic*)]
    \item $\int\frac{\mathrm{d}x}{5-3\cos x}$;
    \item $\int\frac{\mathrm{d}x}{2+\sin^{2}x}$;
    \item $\int\frac{\mathrm{d}x}{1+\tan x}$;
    \item $\int\frac{\tan x}{1+\tan x+\tan^{2}x}\mathrm{d}x$;
    \item $\int\frac{\sin^{3}x}{\cos^{4}x}\mathrm{d}x$;
    \item $\int\frac{\mathrm{d}x}{\cos^{4}x\sin x}$;
    \item $\int\frac{1-\tan x}{1+\tan x}\mathrm{d}x$. (提示: $\frac{1-\tan x}{1+\tan x}=\frac{\cos x-\sin x}{\cos x+\sin x}$ 或 $u=\tan x$.)
  \end{enumerate}
\end{exblock}

% ==================================================================
%  徐森林、薛春华《数学分析》（第一册）
%  第 5 章复习题
% ==================================================================
%  源：原书印刷页 327-328（PDF 物理页 333-334）
%  底本：PaddleOCR-VL-1.6 + 高清影印逐题核验校订
%  规范：exblock 容器; 全西文半角标点; 单一落盘出口 _texout.py
% ==================================================================
\begin{exblock}{复习题 5}

\item 计算不定积分:
\begin{enumerate}[label=(\arabic*)]
  \item $\int\frac{\arcsin x}{x^{2}}\mathrm{d}x$;
  \item $\int\frac{\mathrm{d}x}{\sqrt{\sin x\cos^{7}x}}$;
  \item $\int x\ln\frac{1+x}{1-x}\mathrm{d}x$;
  \item $\int\mathrm{e}^{x}\left(\frac{1-x}{1+x^{2}}\right)^{2}\mathrm{d}x$.
\end{enumerate}

\item 设 $I_{n}=\int\frac{v^{n}}{\sqrt{u}}$, 其中 $u=a_{1}+b_{1}x$, $v=a_{2}+b_{2}x$, 求递推公式.

\item 求不定积分:
\begin{enumerate}[label=(\arabic*)]
  \item $\int\frac{\mathrm{d}x}{(x+a)^{m}(x+b)^{n}}$, $m,n\in\mathbb{N}$;
  \item $\int\frac{\mathrm{d}x}{(1+x^{n})\sqrt[n]{1+x^{n}}}$.
\end{enumerate}

\item 计算不定积分:
\begin{enumerate}[label=(\arabic*)]
  \item $\int\frac{\mathrm{d}x}{2\sin x-\cos x+5}$;
  \item $\int\frac{\cos x\sin x}{\cos x+\sin x}\mathrm{d}x$;
  \item $\int\frac{\sin x}{1+\cos x+\sin x}\mathrm{d}x$;
  \item $\int\frac{\sin^{2}x}{1+\sin^{2}x}\mathrm{d}x$;
  \item $\int\frac{\mathrm{d}x}{\cos^{4}x+\sin^{4}x}$;
  \item $\int\frac{\sin x}{\cos^{3}x+\sin^{3}x}\mathrm{d}x$;
  \item $\int\sqrt{\tan x}\,\mathrm{d}x$;
  \item $\int\sqrt{\tan^{2}x+2}\,\mathrm{d}x$;
  \item $\int\frac{\sin^{2}x\cos x}{\sin x+\cos x}\mathrm{d}x$;
  \item $\int\frac{\mathrm{d}x}{(1+2^{x})^{4}}$;
  \item $\int\sqrt{1+\sin x}\,\mathrm{d}x$;
  \item $\int\frac{x+\sin x}{1+\cos x}\mathrm{d}x$;
  \item $\int\arctan(1+\sqrt{x})\mathrm{d}x$.
\end{enumerate}
\end{exblock}

